%% =====================================================================
%%  Explainable Visual Commonsense Reasoning for Social Interactions
%%  Project Report -- Phase I
%%  Sri Sivasubramaniya Nadar College of Engineering, Dept. of CSE
%%
%%  Compile with:  pdflatex main.tex  (run twice for ToC / references)
%%  Overleaf:      upload this file + a folder "figures/" with
%%                   figures/training_loss.png   (R2 Fig 6.1)
%%                   figures/ui_predictions.png  (R2 Fig 6.3)
%%                   figures/ui_evidence.png     (R2 Fig 6.4)
%%  Missing images are replaced by a red placeholder box, so the file
%%  compiles even before the images are added.
%%  Every \fillin{...} (printed in red) is a value you must supply.
%% =====================================================================
\documentclass[12pt,a4paper,oneside]{report}

%% ---------- Page layout and fonts -----------------------------------
\usepackage[a4paper,left=1.5in,right=1in,top=1in,bottom=1in]{geometry}
\usepackage[T1]{fontenc}
\usepackage[utf8]{inputenc}
\usepackage{newtxtext,newtxmath}
\usepackage{setspace}
\usepackage{microtype}
\setlength{\emergencystretch}{2em}
\onehalfspacing

%% ---------- General packages ----------------------------------------
\usepackage{amsmath}
\usepackage{graphicx}
\usepackage{booktabs}
\usepackage{tabularx}
\usepackage{array}
\usepackage{multirow}
\usepackage{float}
\usepackage[font=small,labelfont=bf]{caption}
\usepackage{enumitem}
\usepackage[table]{xcolor}
\usepackage{titlesec}
\usepackage[nottoc,notlof,notlot]{tocbibind}
\usepackage{tikz}
\usetikzlibrary{positioning,arrows.meta,fit,backgrounds,shapes.geometric,calc}
\usepackage{pgfplots}
\pgfplotsset{compat=1.17}
\usepackage[hidelinks]{hyperref}
\usepackage[chapter]{algorithm}
\usepackage{algpseudocode}

%% ---------- Chapter heading style (CHAPTER n / TITLE) ---------------
\titleformat{\chapter}[display]
  {\normalfont\bfseries\centering}
  {\Large\MakeUppercase{\chaptertitlename}\ \thechapter}
  {1.2ex}
  {\LARGE\MakeUppercase}
\titlespacing*{\chapter}{0pt}{0pt}{28pt}
\titleformat{\section}{\normalfont\large\bfseries}{\thesection}{0.8em}{}
\titleformat{\subsection}{\normalfont\normalsize\bfseries}{\thesubsection}{0.8em}{}
\setcounter{secnumdepth}{3}
\setcounter{tocdepth}{2}

%% ---------- Algorithm style -----------------------------------------
\algrenewcommand\algorithmicrequire{\textbf{Input:}}
\algrenewcommand\algorithmicensure{\textbf{Output:}}

%% ---------- Colours (charts and diagrams) ----------------------------
\definecolor{phaseone}{gray}{0.90}
\definecolor{seriesA}{HTML}{2A78D6}   % blue
\definecolor{seriesB}{HTML}{EB6834}   % orange
\definecolor{seriesC}{HTML}{1BAF7A}   % aqua
\definecolor{neutralbar}{HTML}{9A9994}

%% ---------- Handy macros --------------------------------------------
\newcommand{\fillin}[1]{\textcolor{red}{[#1]}}
\newcommand{\QtoA}{Q$\rightarrow$A}
\newcommand{\QAtoR}{QA$\rightarrow$R}
\newcommand{\QtoAR}{Q$\rightarrow$AR}
\newcommand{\code}[1]{\texttt{#1}}
%% image with automatic placeholder when the file is missing
\newcommand{\figimg}[2]{%
  \IfFileExists{#2}{\includegraphics[width=#1]{#2}}{%
    \fbox{\parbox[c][5.5cm][c]{\dimexpr#1-2\fboxsep-2\fboxrule\relax}{\centering\small
      \textcolor{red}{Place image here:}\\[2pt]\texttt{\detokenize{#2}}}}}}

\renewcommand{\arraystretch}{1.15}
\newcolumntype{L}{>{\raggedright\arraybackslash}X}
\renewcommand{\bibname}{References}

%% ---------- Report details (edit here) ------------------------------
\newcommand{\reporttitle}{Explainable Visual Commonsense Reasoning for Social Interactions}
\newcommand{\guidename}{Dr. M. Saritha}
\newcommand{\guidedesig}{Assistant Professor}
\newcommand{\hodname}{\fillin{Name of the HoD}}
\newcommand{\reportmonth}{\fillin{Month} 2026}

%% =====================================================================
\begin{document}

%% ---------------------------------------------------------------------
%% TITLE PAGE
%% ---------------------------------------------------------------------
\pagenumbering{roman}
\begin{titlepage}
\begin{center}
\vspace*{0.5cm}
{\LARGE\bfseries EXPLAINABLE VISUAL COMMONSENSE\\[4pt]
REASONING FOR SOCIAL INTERACTIONS\par}
\vspace{1.2cm}
{\large\bfseries PROJECT REPORT -- PHASE I\par}
\vspace{1.2cm}
{\large Submitted by\par}
\vspace{0.4cm}
\begin{tabular}{ll}
\textbf{Vijay Sherwin G}   & 3122235001158\\
\textbf{Thirilogsha S S}   & 3122235001147\\
\textbf{Udhaya Sakthi K}   & 3122235001152\\
\end{tabular}

\vspace{1.0cm}
{\large\itshape in partial fulfilment of the requirements for the award of the degree of\par}
\vspace{0.4cm}
{\Large\bfseries BACHELOR OF ENGINEERING\par}
\vspace{0.2cm}
{\large in\par}
\vspace{0.2cm}
{\Large\bfseries COMPUTER SCIENCE AND ENGINEERING\par}
\vspace{1.0cm}
{\large Under the guidance of\par}
\vspace{0.2cm}
{\large\bfseries \guidename\par}
{\large \guidedesig\par}
\vfill
{\large\bfseries Department of Computer Science and Engineering\\[2pt]
Sri Sivasubramaniya Nadar College of Engineering\\[2pt]
Kalavakkam -- 603 110\par}
\vspace{0.4cm}
{\large\bfseries \reportmonth\par}
\end{center}
\end{titlepage}

\setcounter{page}{2}

%% ---------------------------------------------------------------------
%% BONAFIDE CERTIFICATE
%% ---------------------------------------------------------------------
\chapter*{Bonafide Certificate}
\addcontentsline{toc}{chapter}{Bonafide Certificate}

Certified that this project report titled \textbf{``\reporttitle''} is the
bonafide work of \textbf{Vijay Sherwin G (3122235001158)},
\textbf{Thirilogsha S S (3122235001147)} and
\textbf{Udhaya Sakthi K (3122235001152)}, who carried out the Phase I project
work under my supervision.

\vspace{3cm}
\noindent
\begin{minipage}[t]{0.47\textwidth}
\textbf{SIGNATURE}\\[4pt]
\hodname\\
Head of the Department\\
Department of Computer Science and Engineering\\
Sri Sivasubramaniya Nadar College of Engineering
\end{minipage}\hfill
\begin{minipage}[t]{0.47\textwidth}
\textbf{SIGNATURE}\\[4pt]
\guidename\\
Supervisor, \guidedesig\\
Department of Computer Science and Engineering\\
Sri Sivasubramaniya Nadar College of Engineering
\end{minipage}

\vspace{2.5cm}
\noindent Submitted for the project viva-voce examination held on \rule{4cm}{0.4pt}

\vspace{2.5cm}
\noindent
\textbf{INTERNAL EXAMINER} \hfill \textbf{EXTERNAL EXAMINER}

%% ---------------------------------------------------------------------
%% ACKNOWLEDGEMENT
%% ---------------------------------------------------------------------
\chapter*{Acknowledgement}
\addcontentsline{toc}{chapter}{Acknowledgement}

We express our sincere gratitude to our guide, \textbf{\guidename},
\guidedesig, Department of Computer Science and Engineering, for her continuous
guidance, encouragement and valuable feedback throughout the first phase of
this project.

We thank \hodname, Head of the Department of Computer Science and Engineering,
and the members of the project review panel for their suggestions during the
reviews. We are grateful to the management and the Principal of Sri
Sivasubramaniya Nadar College of Engineering for providing the facilities
required for this work.

We also thank the faculty and staff of the Department of Computer Science and
Engineering, and our families and friends, for their support.

\vspace{1.5cm}
\begin{flushright}
Vijay Sherwin G\\
Thirilogsha S S\\
Udhaya Sakthi K
\end{flushright}

%% ---------------------------------------------------------------------
%% ABSTRACT
%% ---------------------------------------------------------------------
\chapter*{Abstract}
\addcontentsline{toc}{chapter}{Abstract}

Understanding a social interaction in an image requires more than recognising
the people and objects in it. The activity taking place, the relationship
between the people and their intention are rarely visible directly; they must
be inferred from visual cues together with everyday commonsense knowledge.

The objective of this project is to build an explainable Visual Commonsense
Reasoning system that predicts the activity, relationship and intention of a
social interaction, and supports every prediction with evidence that a user can
inspect.

The work is carried out in two phases. Phase I, presented in this report,
builds the prediction pipeline. The input is standardised using OpenCV, a
lightweight vision-language model (SmolVLM) describes the people, actions and
setting, ConceptNet supplies related commonsense knowledge, and a reasoning
stage fuses these sources through log-linear fusion to produce the three
predictions with confidence scores, ranked alternatives and a textual
explanation. SmolVLM is also fine-tuned on the VCR dataset with LoRA; the
checkpoint selected from a 20-epoch run (epoch 11) raises accuracy on 300
held-out VCR questions from 44.7\% to 58.7\% for \QtoA, from 41.3\% to 57.7\%
for \QAtoR{} and from 15.3\% to 32.7\% for \QtoAR. The integrated pipeline is
exposed through a FastAPI--React interface and is validated by 23 automated
tests, all of which pass. Phase II will verify the explanations through visual
attribution, cross-modal consistency and faithfulness evaluation, and present
them in an explainability dashboard.

\vspace{0.5cm}
\noindent\textbf{Keywords:} visual commonsense reasoning, social interaction
understanding, vision-language models, SmolVLM, ConceptNet, LoRA,
explainable AI.

%% ---------------------------------------------------------------------
%% LISTS
%% ---------------------------------------------------------------------
\tableofcontents
\listoffigures
\listoftables
\cleardoublepage
\phantomsection
\addcontentsline{toc}{chapter}{List of Algorithms}
\listofalgorithms

\chapter*{List of Abbreviations}
\addcontentsline{toc}{chapter}{List of Abbreviations}
\begin{tabular}{@{}p{3.2cm}p{10.5cm}@{}}
API      & Application Programming Interface\\
CLAHE    & Contrast Limited Adaptive Histogram Equalisation\\
CN       & ConceptNet\\
CUDA     & Compute Unified Device Architecture\\
Grad-CAM & Gradient-weighted Class Activation Mapping\\
HSV      & Hue, Saturation, Value colour space\\
IDF      & Inverse Document Frequency\\
JSON     & JavaScript Object Notation\\
LAB      & CIE $L^*a^*b^*$ colour space\\
LoRA     & Low-Rank Adaptation\\
MPS      & Metal Performance Shaders (Apple GPU backend)\\
PEFT     & Parameter-Efficient Fine-Tuning\\
\QtoA    & Question to Answer (VCR sub-task)\\
\QAtoR   & Question and Answer to Rationale (VCR sub-task)\\
\QtoAR   & Question to Answer and Rationale (joint VCR metric)\\
SDG      & Sustainable Development Goal\\
VCR      & Visual Commonsense Reasoning\\
VLM      & Vision-Language Model\\
VQA      & Visual Question Answering\\
\end{tabular}

\cleardoublepage
\pagenumbering{arabic}

%% =====================================================================
%% CHAPTER 1  INTRODUCTION
%% =====================================================================
\chapter{Introduction}
\label{ch:intro}

Visual understanding has progressed from object recognition to higher-level
reasoning about actions, relationships, context and human intention. Social
interactions are particularly challenging because their interpretation depends
on several people at once, on their posture and gestures, on the objects they
use, on their spatial arrangement and on the wider context. Visual Commonsense
Reasoning (VCR) extends conventional recognition by requiring a model to infer
information that is not explicitly visible in an image~\cite{zellers2019vcr}.
This report presents the first phase of a system that performs such reasoning
for social interactions and exposes the evidence behind each of its
conclusions.

\section{Motivation}

A photograph of two people shaking hands may show a business deal, a greeting
between friends, a sporting gesture or a reconciliation after an argument. The
pixels are much the same in each case; what differs is the context and the
commonsense knowledge a person brings to the scene. Humans make these
inferences effortlessly, but automated systems that make them are increasingly
used in settings where the inference matters: assistive tools that describe
scenes to visually impaired users, social robots that must decide how to
approach people, and tools that summarise or moderate visual media.

Two problems limit the use of current models in such settings.

\begin{enumerate}[leftmargin=*]
\item \textbf{Opacity.} High-accuracy VCR models select an answer from a set of
  options but do not show which visual cues or which pieces of world knowledge
  led to the choice. Large generative vision-language models do produce
  explanations, but these explanations are not guaranteed to reflect the
  evidence the model actually used, and they can describe details that are not
  present in the image.
\item \textbf{Cost.} The strongest multimodal models have tens of billions of
  parameters and require data-centre hardware, which rules them out for
  on-device, privacy-preserving or low-cost deployment.
\end{enumerate}

This project is motivated by the need for a system that is both
\emph{lightweight} and \emph{inspectable}: it should run on modest hardware, and
every prediction about a social interaction should come with structured
evidence -- the visual concepts detected, the commonsense facts used, the
probabilities from each source and the alternatives that were considered --
that a user can check.

\section{Background}

\subsection{Visual Commonsense Reasoning}

The VCR task, introduced by Zellers et al.~\cite{zellers2019vcr}, poses a
multiple-choice question about an image taken from a movie scene. The people
and objects in the image are referred to by tags such as \emph{person1} and
\emph{person2}. A model must solve two sub-tasks: \QtoA, choosing the correct
answer from four options, and \QAtoR, choosing the rationale that justifies
the correct answer from four options. The joint metric \QtoAR{} counts a
question as correct only when both the answer and the rationale are correct.
With four options each, chance accuracy is 25\% for \QtoA{} and \QAtoR{} and
6.25\% for \QtoAR.

\subsection{Vision-Language Models}

A vision-language model (VLM) couples an image encoder with a language model so
that it can describe images and answer questions about them. Modern VLMs encode
the image into a sequence of visual tokens, project them into the embedding
space of a pre-trained language model, and generate text
autoregressively~\cite{liu2023llava,li2023blip2}. SmolVLM~\cite{marafioti2025smolvlm}
is a family of compact VLMs designed for efficient inference. The variant used
in this project, \code{SmolVLM-Instruct}, has about 2.2 billion parameters and
combines a SigLIP vision encoder~\cite{zhai2023siglip} with a small
language model through a connector that compresses the visual tokens.

\subsection{Commonsense Knowledge Graphs}

ConceptNet~\cite{speer2017conceptnet} is a multilingual knowledge graph of
everyday knowledge. Its nodes are words and phrases (for example
\emph{concert}, \emph{sing}, \emph{stage}) and its edges are labelled
relations such as \emph{RelatedTo}, \emph{IsA}, \emph{UsedFor},
\emph{AtLocation}, \emph{CapableOf}, \emph{Antonym} and \emph{DistinctFrom},
each with a confidence weight. Because each edge is a readable statement
(``a performer is typically found at a stage''), ConceptNet is well suited to
producing evidence that a human can inspect.

\subsection{Parameter-Efficient Fine-Tuning}

Fine-tuning all the weights of a multi-billion-parameter model is expensive.
Low-Rank Adaptation (LoRA)~\cite{hu2022lora} freezes the pre-trained weight
matrix $W_0 \in \mathbb{R}^{d \times k}$ and learns a low-rank update:
\begin{equation}
  h = W_0 x + \frac{\alpha}{r} B A x, \qquad
  B \in \mathbb{R}^{d \times r},\; A \in \mathbb{R}^{r \times k},\; r \ll \min(d,k)
  \label{eq:lora}
\end{equation}
where $r$ is the rank and $\alpha$ a scaling factor. Only $A$ and $B$ are
trained, so the number of trainable parameters and the optimiser memory are a
small fraction of full fine-tuning, and after training $BA$ can be merged into
$W_0$ so that inference is no slower than with the base model.

\subsection{Explainability and Faithfulness}

An explanation is \emph{plausible} if it convinces a human and \emph{faithful}
if it reflects the reasoning the model actually performed. Visual explanation
methods such as Grad-CAM~\cite{selvaraju2017gradcam} highlight the image
regions that most influence a prediction. Their faithfulness can be measured
with the Insertion--Deletion metric~\cite{petsiuk2018rise}, which tracks how the
model's confidence changes as the highlighted pixels are added or removed, and
with the Pointing Game~\cite{zhang2018excitation}, which checks whether the most
salient point falls inside the relevant object. Textual explanations can be
compared with visual explanations to measure cross-modal
consistency~\cite{zhang2024crossmodal,parcalabescu2025faithfulness}. These
methods form the basis of Phase II of this project.

\section{Problem Definition}

A social interaction may be visually ambiguous even when the people and objects
in the scene are detected correctly. A system may recognise two people and
their gestures but still need commonsense knowledge to distinguish between
possible activities, or to infer the relationship and intention behind the
interaction. The problem addressed in this project is therefore:

\begin{quote}
\itshape Given a static image, video frame, or live-camera frame containing zero
or more people, develop a lightweight system that identifies the visual
concepts in the scene, enriches them with commonsense knowledge, and predicts
the activity, relationship, and intention of a social interaction with
interpretable evidence.
\end{quote}

\paragraph{Inputs.} The system accepts a static image, supplied as a file path,
a URL, a NumPy array or a PIL image, as well as sampled frames from a video file
and frames from a live camera. Along with the visual input, the user may ask an
optional free-form question about the scene and may also supply
multiple-choice answer options for that question.

\paragraph{Outputs.} The principal output is a structured prediction of the
activity, relationship and intention of the interaction. Each of the three
predictions carries a label, a confidence value and a ranked list of
alternative labels, and every alternative records its VLM probability,
commonsense probability and draft-match value. The prediction is accompanied by
its supporting evidence and a generated textual explanation. The intermediate
results are returned as well: the visual concepts and draft reasoning produced
by the VLM, the ConceptNet facts and label-support scores, and the processing
time of each module.

\paragraph{Constraints.} The system is built around lightweight multimodal
models so that it can run on modest hardware. The ConceptNet integration uses a
local SQLite database for reproducible offline operation and can fall back to
other configured backends. The label space used for reasoning is configurable
through \path{configs/taxonomy.yaml}.

\paragraph{Scope of Phase I.} The complete framework follows the workflow
\begin{center}
\small\textbf{Input $\rightarrow$ Preprocessing $\rightarrow$ VLM $\rightarrow$
Commonsense $\rightarrow$ Reasoning $\rightarrow$ Explanation $\rightarrow$
Consistency $\rightarrow$ Evaluation $\rightarrow$ Dashboard}
\end{center}
Phase I, reported here, covers the first four modules -- Input and
Preprocessing, the Small Vision-Language Model, ConceptNet Commonsense
Enrichment and Social Interaction Reasoning -- together with the offline
fine-tuning of SmolVLM on VCR. Its deliverable is a working pipeline that takes
visual input and produces a social-interaction prediction with a draft
explanation. Phase II adds the textual and visual explanation paths,
cross-modal consistency analysis, explainability evaluation and the final
explainability dashboard.

\section{Organization of the Report}

The rest of this report is organised as follows.
Chapter~\ref{ch:lit} surveys the literature on visual commonsense reasoning,
vision-language models, knowledge-augmented reasoning, social relation
recognition and explainability, identifies the gaps this project addresses and
states the research objectives.
Chapter~\ref{ch:method} describes the proposed methodology: the overall
architecture, each Phase I module with its algorithm, the LoRA fine-tuning
procedure and the integration of the modules.
Chapter~\ref{ch:results} describes the datasets and the hardware and software
ecosystem, the experiments conducted, and analyses the results.
Chapter~\ref{ch:impact} discusses the social impact and sustainability of the
work. Chapter~\ref{ch:conclusion} concludes the report and outlines the work
planned for Phase II. The appendices list the configuration parameters and the
API endpoints of the implementation.

%% =====================================================================
%% CHAPTER 2  LITERATURE SURVEY
%% =====================================================================
\chapter{Literature Survey}
\label{ch:lit}

This chapter reviews work in five areas that bear on the problem: models and
datasets for visual commonsense reasoning, pre-trained and instruction-tuned
vision-language models, reasoning with external knowledge, recognition of
social relations, and the explanation and evaluation of multimodal models.

\section{Papers Related to the Problem Identified}

\subsection{Visual Commonsense Reasoning Models}

\textbf{Zellers et al.}~\cite{zellers2019vcr} introduced the VCR dataset and the
Recognition-to-Cognition network (R2C). R2C grounds the words of the question
and the candidate responses in detected object regions, contextualises the
response with respect to the question and the image, and reasons over the
result to score each option. The dataset was constructed with
\emph{Adversarial Matching}, which reuses correct answers from other questions
as distractors so that the wrong options are fluent and relevant. The paper
showed a large gap between the best model and human performance, establishing
VCR as a benchmark for cognition-level understanding.

\textbf{Lin et al.}~\cite{lin2019tabvcr} proposed TAB-VCR, a simpler baseline
that attaches attributes to the detected objects and adds new tags for objects
that are mentioned in the text but not tagged in the image. The improved visual
grounding gave higher accuracy than R2C with fewer parameters, showing that
better grounding of text in image regions matters more than model complexity.

\textbf{Yu et al.}~\cite{yu2019hgl} proposed Heterogeneous Graph Learning (HGL),
which builds a vision-to-answer graph and a question-to-answer graph and uses a
contextual voting module to exploit long-range visual context. It models the
relations between visual and linguistic elements explicitly, but its reasoning
remains internal to the network.

\textbf{Park et al.}~\cite{park2020visualcomet} introduced VisualCOMET, a
dataset of more than a million textual commonsense inferences about still
images, describing what may have happened before, what may happen next and the
intents of the people in the scene. It frames visual commonsense as the
generation of inferences rather than the selection of options, which is closer
to the explanations this project produces.

\subsection{Pre-trained Vision-Language Transformers}

\textbf{ViLBERT}~\cite{lu2019vilbert} extended BERT to a two-stream architecture
in which visual and textual streams interact through co-attentional transformer
layers, pre-trained on image--caption pairs and transferred to tasks including
VCR. \textbf{VL-BERT}~\cite{su2020vlbert} used a single-stream transformer over
both modalities, and \textbf{UNITER}~\cite{chen2020uniter} learned a universal
image--text representation from several large datasets with word--region
alignment objectives, reaching state-of-the-art VCR accuracy at the time.
\textbf{MERLOT}~\cite{zellers2021merlot} learned multimodal script knowledge from
millions of YouTube videos and their transcripts, and transferred well to VCR.

These models raised VCR accuracy considerably, but they are discriminative
classifiers: they score given options and provide no explanation of their own
beyond attention weights, which are not reliable indicators of the evidence
used.

\subsection{Large and Small Vision-Language Models}

\textbf{CLIP}~\cite{radford2021clip} trained image and text encoders
contrastively on 400 million image--text pairs, enabling zero-shot recognition
from natural-language labels. \textbf{BLIP-2}~\cite{li2023blip2} connected a
frozen image encoder to a frozen large language model through a lightweight
Querying Transformer, and \textbf{LLaVA}~\cite{liu2023llava} connected a CLIP
encoder to an instruction-tuned language model and trained it on
machine-generated visual instruction data. These generative VLMs answer open
questions and produce fluent free-form explanations, but the largest of them are
expensive to run and their explanations can include details that are not
present in the image.

\textbf{SmolVLM}~\cite{marafioti2025smolvlm} targets the opposite end of the
design space: compact models with aggressive visual-token compression, intended
for on-device and low-memory inference. Its vision encoder is based on
SigLIP~\cite{zhai2023siglip}. SmolVLM makes a VLM practical on a laptop, which
is why it is used in this project, but on its own it is not adapted to
commonsense reasoning about people.

\subsection{Knowledge-Augmented Visual Reasoning}

\textbf{ConceptNet 5.5}~\cite{speer2017conceptnet} provides a large graph of
general knowledge whose edges are readable statements.
\textbf{OK-VQA}~\cite{marino2019okvqa} is a VQA benchmark whose questions cannot
be answered from the image alone and require external knowledge.
\textbf{KRISP}~\cite{marino2021krisp} combined the implicit knowledge of a
pre-trained transformer with symbolic knowledge from graphs, including
ConceptNet, through a graph network, and showed that both kinds of knowledge
contribute. In these systems, however, the retrieved knowledge is usually
converted into hidden features, so the facts that influenced an answer are not
shown to the user.

\subsection{Social Relation and Interaction Recognition}

\textbf{Sun et al.}~\cite{sun2017domain} annotated social relations in the PIPA
photo-album dataset following a domain-based theory from social psychology and
predicted relations from attributes of the people such as age, gender, clothing
and proximity. \textbf{Li et al.}~\cite{li2017dualglance} introduced the PISC
dataset and a Dual-Glance model that first looks at the pair of people and then
attends to contextual objects in the scene. These works show that context and
person attributes determine the relationship, but they treat the relationship
as a single classification target and do not reason jointly about the activity
and the intention.

\subsection{Parameter-Efficient Adaptation}

\textbf{LoRA}~\cite{hu2022lora} adapts large models by training low-rank updates
to frozen weight matrices, and \textbf{QLoRA}~\cite{dettmers2023qlora} combines
LoRA with 4-bit quantisation of the frozen weights so that large models can be
fine-tuned on a single GPU. These methods make it feasible to adapt a VLM to VCR
on free cloud notebooks.

\subsection{Explainability and Faithfulness Evaluation}

\textbf{Grad-CAM}~\cite{selvaraju2017gradcam} produces class-specific saliency
maps from the gradients flowing into the last convolutional or attention layer.
\textbf{RISE}~\cite{petsiuk2018rise} explains black-box models by random masking
and introduced the Insertion and Deletion metrics for evaluating saliency maps.
The \textbf{Pointing Game}~\cite{zhang2018excitation} checks whether the
maximum of a saliency map lies on the annotated object.
\textbf{e-ViL}~\cite{kayser2021evil} provides a benchmark for natural-language
explanations in vision-language tasks. Recent work examines whether the textual
explanations of VLMs are self-consistent and faithful to their use of the image
and the text~\cite{parcalabescu2025faithfulness}, and whether multimodal models
give consistent answers across modalities~\cite{zhang2024crossmodal}.

\section{Summary of the Literature}

Table~\ref{tab:litsummary} summarises the principal works and their limitations
with respect to the problem defined in Chapter~\ref{ch:intro}.

\begin{table}[!htb]
\centering
\caption{Summary of related work}
\label{tab:litsummary}
\footnotesize
\renewcommand{\arraystretch}{1.25}
\begin{tabularx}{\textwidth}{@{}p{2.6cm}p{3.6cm}L@{}}
\toprule
\textbf{Work} & \textbf{Approach} & \textbf{Limitation for this problem}\\
\midrule
R2C~\cite{zellers2019vcr} & Grounded attention over detections; VCR benchmark & Selects options; no evidence shown to the user \\
TAB-VCR~\cite{lin2019tabvcr} & Attributes and new tags for better grounding & Same discriminative setting; no external knowledge \\
HGL~\cite{yu2019hgl} & Heterogeneous vision/question--answer graphs & Graph reasoning is internal and not readable \\
ViLBERT, VL-BERT, UNITER~\cite{lu2019vilbert,su2020vlbert,chen2020uniter} & Large-scale vision--language pre-training & High accuracy but opaque; heavy fine-tuning \\
MERLOT~\cite{zellers2021merlot} & Script knowledge from video & Very large pre-training cost \\
BLIP-2, LLaVA~\cite{li2023blip2,liu2023llava} & Generative VLMs with LLM backbones & Fluent but unverified explanations; large models \\
SmolVLM~\cite{marafioti2025smolvlm} & Compact efficient VLM & Not adapted to social commonsense reasoning \\
KRISP~\cite{marino2021krisp} & Implicit + symbolic (ConceptNet) knowledge & Knowledge hidden in features; not social-specific \\
Dual-Glance, PIPA~\cite{li2017dualglance,sun2017domain} & Social relation classification with context & Relationship only; no activity or intention \\
Grad-CAM, RISE, Pointing Game~\cite{selvaraju2017gradcam,petsiuk2018rise,zhang2018excitation} & Saliency and its evaluation & Not linked to textual explanations \\
e-ViL~\cite{kayser2021evil} & Benchmark for natural-language explanations & Evaluates text alone, not cross-modal agreement \\
\bottomrule
\end{tabularx}
\end{table}

\section{Gaps Identified}

The survey reveals the following gaps.

\begin{enumerate}[label=\textbf{G\arabic*.},leftmargin=*]
\item \textbf{Accuracy without explanation.} The strongest VCR models choose an
  answer and a rationale from given options but do not expose the visual cues
  or knowledge behind the choice.
\item \textbf{Explanations without verification.} Generative VLMs produce
  fluent explanations, but these are not checked against the visual evidence
  and can describe things that are not in the image.
\item \textbf{Hidden knowledge.} Knowledge-augmented models fold the retrieved
  facts into hidden representations, so the user cannot see which facts
  influenced a prediction.
\item \textbf{Single-target social understanding.} Social relation recognition
  treats the relationship as an isolated label; the activity, the relationship
  and the intention are not predicted jointly and consistently.
\item \textbf{Separate evaluation of explanations.} Visual saliency and textual
  explanations are evaluated separately; their mutual consistency is rarely
  measured.
\item \textbf{Computational cost.} Most high-performing systems require large
  models and data-centre hardware, which prevents lightweight or on-device use.
\end{enumerate}

\section{Research Objectives}
\label{sec:objectives}

To address these gaps, the project sets the following objectives. Objectives
1--10 are the Phase I objectives addressed in this report; objectives 11--13
belong to Phase II.

\begin{enumerate}[leftmargin=*]
\item To implement a robust input and preprocessing pipeline for images, video
  frames and live-camera input using OpenCV.
\item To integrate a lightweight vision-language model, SmolVLM, for visual
  understanding and language-based reasoning (G6).
\item To fine-tune SmolVLM on the VCR dataset using LoRA and select the best
  checkpoint on held-out validation questions (G6).
\item To extract people, objects, actions, gestures, setting and spatial
  information from the visual input as structured concepts (G1).
\item To enrich the visual concepts with everyday commonsense knowledge from
  ConceptNet and expose that knowledge as readable evidence paths (G3).
\item To develop a social-interaction reasoning module that fuses visual,
  commonsense and VLM draft evidence (G1, G3).
\item To predict the activity, relationship and intention of a social
  interaction in sequence so that they are mutually consistent (G4).
\item To provide confidence values, ranked alternatives and supporting evidence
  for each prediction (G1).
\item To integrate the Phase I modules into a reusable Python pipeline and an
  operational FastAPI--React demonstration interface.
\item To validate the implemented modules using automated unit and end-to-end
  tests.
\item \emph{(Phase II)} To generate independent textual and visual (Grad-CAM)
  explanations for each prediction (G2).
\item \emph{(Phase II)} To measure the cross-modal consistency of the two
  explanations and evaluate their faithfulness with the Insertion--Deletion
  metric and the Pointing Game (G2, G5).
\item \emph{(Phase II)} To present the predictions, explanations and scores in
  an explainability dashboard.
\end{enumerate}

%% =====================================================================
%% CHAPTER 3  PROPOSED METHODOLOGY
%% =====================================================================
\chapter{Proposed Methodology}
\label{ch:method}

This chapter describes the architecture of the proposed system and the design
of each Phase I module. The Phase I pipeline is
\begin{center}
\textbf{Input $\rightarrow$ OpenCV Preprocessing $\rightarrow$ SmolVLM
$\rightarrow$ ConceptNet $\rightarrow$ Social Reasoning}
\end{center}
with an offline training phase in which SmolVLM is fine-tuned on VCR.

\section{Architectural Design of the Proposed System}

The overall architecture, covering both phases, is shown in
Figure~\ref{fig:overall}. Visual input is preprocessed with OpenCV and passed to
a small vision-language model, which can be replaced by a version fine-tuned
offline on the VCR dataset. The output of the model is enriched with
commonsense knowledge, and the social interaction reasoning module fuses the two
into a prediction.

In Phase II, this prediction is explained along two independent paths, one
textual and one visual. The consistency between the two explanations is
measured, their faithfulness is evaluated with the Insertion--Deletion metric
and the Pointing Game, and the results are presented in an explainability
dashboard.

\begin{figure}[p]
\centering
\resizebox{\linewidth}{!}{%
\begin{tikzpicture}[
  font=\small, >={Stealth[length=2mm]},
  box/.style={draw, rounded corners=2pt, fill=white, align=center,
              minimum height=0.75cm, inner sep=3pt},
  grp/.style={draw, rounded corners=4pt, inner sep=0.18cm},
  p1/.style={grp, fill=phaseone},
  p2/.style={grp, fill=white},
  ttl/.style={anchor=west, font=\small\bfseries, inner sep=0pt},
  arr/.style={->, thick}
]
%% ---- Input and preprocessing (Phase I)
\node[ttl] (t1) at (0.6,0.75) {Input and Preprocessing};
\node[box, minimum width=4.2cm] (in1) at (2.9,0.0) {Static image};
\node[box, minimum width=4.2cm] (in2) at (2.9,-0.95) {Video stream / live camera};
\node[box, minimum width=5cm, minimum height=1.3cm] (pre) at (11,-0.475) {OpenCV preprocessing};
\coordinate (c1a) at (0.6,0.75); \coordinate (c1b) at (15.4,-1.33);
\begin{scope}[on background layer]
  \node[p1, fit=(t1)(in1)(in2)(pre)(c1a)(c1b)] (g1) {};
\end{scope}
\draw[arr] (in1.east) -- (in1.east -| pre.west);
\draw[arr] (in2.east) -- (in2.east -| pre.west);

%% ---- Offline training phase
\node[ttl] (t2) at (0.6,-2.35) {Training Phase (Offline)};
\node[box, minimum width=3.4cm] (tr1) at (2.7,-3.15) {VCR dataset};
\node[box, minimum width=3.4cm] (tr2) at (2.7,-4.15) {Data preparation};
\node[box, minimum width=3.4cm] (tr3) at (2.7,-5.15) {LoRA fine-tuning};
\node[box, minimum width=3.4cm] (tr4) at (2.7,-6.15) {Fine-tuned SmolVLM};
\coordinate (c2a) at (0.6,-2.35); \coordinate (c2b) at (4.8,-6.55);
\begin{scope}[on background layer]
  \node[grp, dashed, fill=phaseone!50, fit=(t2)(tr1)(tr4)(c2a)(c2b)] (g2) {};
\end{scope}
\draw[arr] (tr1) -- (tr2);
\draw[arr] (tr2) -- (tr3);
\draw[arr] (tr3) -- (tr4);

%% ---- SmolVLM
\node[ttl] (t3) at (5.6,-2.35) {Small Vision-Language Model (SmolVLM)};
\node[box, minimum width=4.2cm] (v1) at (8.0,-3.5) {Visual understanding};
\node[box, minimum width=4.2cm] (v2) at (13.1,-3.5) {Language and reasoning};
\coordinate (c3a) at (5.6,-2.35); \coordinate (c3b) at (15.4,-4.1);
\begin{scope}[on background layer]
  \node[p1, fit=(t3)(v1)(v2)(c3a)(c3b)] (g3) {};
\end{scope}
\draw[arr] (pre.south) -- (pre.south |- g3.north);
\draw[arr] (v1) -- (v2);
\draw[arr] (tr4.east) -- ++(0.75,0) |- (g3.west);

%% ---- ConceptNet
\node[box, fill=phaseone, font=\small\bfseries, minimum width=9.8cm,
      minimum height=0.85cm] (cn) at (10.5,-6.15)
      {Commonsense Knowledge Enrichment (ConceptNet)};
\draw[arr] (v2.south) -- (v2.south |- cn.north);

%% ---- Social interaction reasoning
\node[ttl] (t4) at (0.6,-7.35) {Social Interaction Reasoning};
\node[box, minimum width=2.5cm, minimum height=1.0cm] (s1) at (1.85,-8.35)  {Context\\fusion};
\node[box, minimum width=2.5cm, minimum height=1.0cm] (s2) at (4.925,-8.35) {Sequential\\prediction};
\node[box, minimum width=2.5cm, minimum height=1.0cm] (s3) at (8.0,-8.35)   {Candidate\\shortlisting};
\node[box, minimum width=2.5cm, minimum height=1.0cm] (s4) at (11.075,-8.35){Log-linear\\fusion};
\node[box, minimum width=2.5cm, minimum height=1.0cm] (s5) at (14.15,-8.35) {Context /\\query};
\coordinate (c4a) at (0.6,-7.35); \coordinate (c4b) at (15.4,-8.85);
\begin{scope}[on background layer]
  \node[p1, fit=(t4)(s1)(s5)(c4a)(c4b)] (g4) {};
\end{scope}
\draw[arr] (s1) -- (s2);
\draw[arr] (s2) -- (s3);
\draw[arr] (s3) -- (s4);
\draw[arr] (s5) -- (s4);
\draw[arr] (cn.south) -- (cn.south |- g4.north);

%% ---- Phase II: explanation paths and consistency
\node[box, font=\small\bfseries, minimum width=7cm, minimum height=0.85cm] (tex) at (4.2,-10.2) {Textual explanation path};
\node[box, font=\small\bfseries, minimum width=7cm, minimum height=0.85cm] (vis) at (11.8,-10.2) {Visual explanation path};
\draw[arr] (tex.north |- g4.south) -- (tex.north);
\draw[arr] (vis.north |- g4.south) -- (vis.north);
\node[box, font=\small\bfseries, minimum width=5.4cm, minimum height=0.85cm] (cons) at (8,-11.5) {Cross-modal consistency analysis};
\draw[arr] (tex.south) |- (cons.west);
\draw[arr] (vis.south) |- (cons.east);

%% ---- Phase II: evaluation
\node[ttl] (t5) at (0.6,-12.55) {Explainability Evaluation};
\node[box, minimum width=6.8cm] (e1) at (4.2,-13.4)  {Insertion--Deletion metric};
\node[box, minimum width=6.8cm] (e2) at (11.8,-13.4) {Pointing game};
\coordinate (c5a) at (0.6,-12.55); \coordinate (c5b) at (15.4,-13.85);
\begin{scope}[on background layer]
  \node[p2, fit=(t5)(e1)(e2)(c5a)(c5b)] (g5) {};
\end{scope}
\draw[arr] (cons.south) -- (cons.south |- g5.north);

%% ---- Phase II: dashboard
\node[ttl] (t6) at (0.6,-14.95) {Final Explainable Output (Dashboard)};
\node[box, font=\footnotesize, minimum width=3.4cm, minimum height=1.0cm] (d1) at (2.3,-15.85)  {Input frame};
\node[box, font=\footnotesize, minimum width=3.4cm, minimum height=1.0cm] (d2) at (6.1,-15.85)  {Predicted\\interaction};
\node[box, font=\footnotesize, minimum width=3.4cm, minimum height=1.0cm] (d3) at (9.9,-15.85)  {Activity, relationship,\\intention};
\node[box, font=\footnotesize, minimum width=3.4cm, minimum height=1.0cm] (d4) at (13.7,-15.85) {Consistency\\score};
\coordinate (c6a) at (0.6,-14.95); \coordinate (c6b) at (15.4,-16.35);
\begin{scope}[on background layer]
  \node[p2, fit=(t6)(d1)(d4)(c6a)(c6b)] (g6) {};
\end{scope}
\draw[arr] (cons.south |- g5.south) -- (cons.south |- g6.north);
\end{tikzpicture}%
}
\caption[Overall architecture of the proposed system]{Overall architecture of
the proposed system (shaded blocks: Phase I, implemented in this report;
unshaded blocks: Phase II)}
\label{fig:overall}
\end{figure}

\section{Phase I Architecture}

The part of the architecture implemented in Phase I is shown in
Figure~\ref{fig:phase1}. The modules are connected through structured Python
data objects rather than by passing unstructured text between stages. The main
implementation files are listed in Table~\ref{tab:codemap}.

\begin{figure}[htbp]
\centering
\begin{tikzpicture}[
  font=\small, >={Stealth[length=2mm]},
  box/.style={draw, rounded corners=2pt, align=center, minimum height=0.8cm,
              minimum width=6.4cm, inner sep=3pt},
  mod/.style={box, fill=phaseone, font=\small\bfseries},
  arr/.style={->, thick}
]
\node[box] (a0) at (0,0) {Image / Video frame / Live camera};
\node[mod] (a1) at (0,-1.2) {Input and Preprocessing};
\node[mod] (a2) at (0,-2.4) {Small Vision-Language Model};
\node[mod] (a3) at (0,-3.6) {ConceptNet Commonsense Enrichment};
\node[mod] (a4) at (0,-4.8) {Social Interaction Reasoning};
\node[box] (a5) at (0,-6.0) {Activity $\cdot$ Relationship $\cdot$ Intention};
\draw[arr] (a0) -- (a1);
\draw[arr] (a1) -- (a2);
\draw[arr] (a2) -- (a3);
\draw[arr] (a3) -- (a4);
\draw[arr] (a4) -- (a5);
\node[draw, cylinder, shape border rotate=90, aspect=0.25, minimum height=1.3cm,
      minimum width=1.9cm, align=center, font=\footnotesize] (db) at (-5.0,-3.6)
      {ConceptNet\\(SQLite)};
\draw[arr] (db.east) -- (a3.west);
\node[draw, dashed, rounded corners=2pt, align=center, font=\footnotesize,
      minimum height=0.8cm] (lora) at (5.1,-2.4) {LoRA adapter\\(VCR, epoch 11)};
\draw[arr] (lora.west) -- (a2.east);
\end{tikzpicture}
\caption{Phase I architecture}
\label{fig:phase1}
\end{figure}

\begin{table}[htbp]
\centering
\caption{Phase I module-to-code mapping}
\label{tab:codemap}
\small
\begin{tabularx}{\textwidth}{@{}p{5.2cm}L@{}}
\toprule
\textbf{Module} & \textbf{Implementation}\\
\midrule
Input and Preprocessing & \path{svcr/preprocessing.py}, \code{Preprocessor}, \code{FrameStream}\\
Small Vision-Language Model & \path{svcr/vlm.py}, \code{SmolVLMEngine}\\
ConceptNet Commonsense Enrichment & \path{svcr/conceptnet/client.py}, \path{enrichment.py}, \path{relations.py}\\
Social Interaction Reasoning & \path{svcr/reasoning.py}, \code{SocialInteractionReasoner}\\
Fine-Tuning (offline) & \path{finetune/train_vcr_lora.py}, \path{configs/finetuned.yaml}\\
Pipeline Integration & \path{svcr/pipeline.py}, \code{Phase1Pipeline}\\
Data Structures & \path{svcr/schemas.py}\\
Configuration & \path{configs/default.yaml}, \path{configs/taxonomy.yaml}\\
\bottomrule
\end{tabularx}
\end{table}

\section{Data Flow}

For a single image, data flows through the pipeline as follows. The image is
acquired and its quality is assessed; the necessary preprocessing is applied,
and the frame is converted into the RGB representation expected by the model
together with a normalised tensor. SmolVLM then examines the frame to identify
the people, objects, actions, gestures, setting and spatial relationships, and
produces a draft activity, relationship, intention and rationale, as well as an
answer to the user's question when one is asked.

Salient concepts are next extracted from these outputs and looked up in
ConceptNet. The useful commonsense relations are retained and used to calculate
the support for each candidate label in the taxonomy. Finally, the VLM
probabilities, the ConceptNet support and the draft-matching scores are fused
to predict the activity, relationship and intention in sequence, and the
structured result is returned with its evidence, alternatives, explanation and
timings.

\section{Input and Preprocessing}
\label{sec:preproc}

The Input and Preprocessing module provides a common input representation for
all downstream components. It is implemented in \path{svcr/preprocessing.py}
using OpenCV~\cite{bradski2000opencv}, NumPy and, where required, PIL.

\paragraph{Input sources.} The module accepts image file paths, HTTP and HTTPS
image URLs, NumPy arrays and PIL images, as well as video files and live camera
streams. Each is converted into a common BGR frame before further processing.

\paragraph{Quality assessment.} For every frame the brightness, contrast and
blur are calculated; blur is estimated from the variance of the Laplacian.
Frames that are dark or low in contrast are enhanced by applying
CLAHE~\cite{zuiderveld1994clahe} to the luminance channel of the LAB colour
space, followed by a mild gamma correction. Optional non-local-means denoising
is available through the configuration.

\paragraph{Resize and normalisation.} The default resizing mode preserves the
aspect ratio and limits the longest side to a size compatible with the model.
The RGB image is retained as an 8-bit image for the SmolVLM processor, and a
normalised tensor is also generated:
\begin{equation}
  x_{\text{norm}} = \frac{x/255 - \mu}{\sigma}, \qquad
  \mu = \sigma = (0.5, 0.5, 0.5)
  \label{eq:norm}
\end{equation}
The tensor has shape $1 \times 3 \times H \times W$ and is retained for the
Phase II operations, Grad-CAM and Insertion--Deletion. The procedure is given in
Algorithm~\ref{alg:preprocess}.

\begin{algorithm}[htbp]
\caption{Input preprocessing}
\label{alg:preprocess}
\begin{algorithmic}[1]
\Require Image source $S$ (path, URL, array or PIL image); target longest side $N$; normalisation constants $\mu,\sigma$
\Ensure Preprocessed frame $F$
\State $I \gets \textsc{LoadAsBGR}(S)$;\; $(h,w) \gets$ size of $I$
\State $g \gets$ greyscale version of $I$
\State $b \gets \operatorname{mean}(g)$;\; $c \gets \operatorname{std}(g)$;\; $z \gets \operatorname{Var}(\nabla^2 g)$ \Comment{brightness, contrast, blur}
\If{$b < 60$ \textbf{or} $c < 25$}
  \State $I \gets$ CLAHE on the L channel of LAB, then gamma $0.8$
\EndIf
\If{denoising is enabled}
  \State $I \gets \textsc{NonLocalMeansDenoise}(I)$
\EndIf
\State $s \gets N / \max(h,w)$
\If{$s < 1$}
  \State $I \gets \textsc{Resize}(I, s)$ \Comment{aspect ratio preserved}
\EndIf
\State $R \gets \textsc{BGRtoRGB}(I)$
\State $T \gets (R/255 - \mu)/\sigma$, arranged as $1 \times 3 \times H \times W$
\State \Return $F = \big(R, T, (w,h), s, \{b, c, z, \text{blurry} = z < 60\}\big)$
\end{algorithmic}
\end{algorithm}

\paragraph{Video and live-camera handling.} For video and live-camera input,
\code{FrameStream} samples frames at a configurable interval. Blurry frames can
be skipped, and frames that show the same scene as the last analysed frame are
filtered out using the Bhattacharyya distance between HSV histograms, with a
periodic refresh so that a scene is not suppressed indefinitely
(Algorithm~\ref{alg:framesample}). For uploaded videos, the web interface
selects up to seven key moments: the clip is divided into scene segments using
colour histograms, and the sharpest frame from the central part of each segment
is analysed.

\begin{algorithm}[htbp]
\caption{Frame sampling for video and live camera}
\label{alg:framesample}
\begin{algorithmic}[1]
\Require Frame stream $\{(f_t, t)\}$; interval $\Delta$; scene-change threshold $\tau$
\Ensure Sequence of frames to analyse
\State $t_{\text{check}} \gets -\infty$;\; $t_{\text{last}} \gets -\infty$;\; $H_{\text{last}} \gets \textsc{null}$
\ForAll{frames $(f_t, t)$ in the stream}
  \If{$t - t_{\text{check}} < \Delta$ \textbf{or} $f_t$ is blurry}
    \State \textbf{continue}
  \EndIf
  \State $t_{\text{check}} \gets t$;\; $H \gets$ HSV histogram of $f_t$
  \If{$H_{\text{last}} \neq \textsc{null}$ \textbf{and} $t - t_{\text{last}} < 3\Delta$ \textbf{and} $d_{\text{Bhatt}}(H_{\text{last}}, H) < \tau$}
    \State \textbf{continue} \Comment{same scene as last analysed frame}
  \EndIf
  \State $t_{\text{last}} \gets t$;\; $H_{\text{last}} \gets H$
  \State \textbf{output} $f_t$
\EndFor
\end{algorithmic}
\end{algorithm}

\section{Small Vision-Language Model}
\label{sec:vlm}

The VLM module is implemented by \code{SmolVLMEngine} in \path{svcr/vlm.py}.
The default configuration uses \code{HuggingFaceTB/SmolVLM-Instruct}, loaded
through Hugging Face Transformers~\cite{wolf2020transformers}. The engine loads
the model and processor once, selects CUDA, MPS or the CPU as available, and
provides three operations: \code{generate()} for deterministic text generation,
\code{score\_options()} for multiple-choice probability estimation, and
\code{run()}, which performs the complete visual-understanding and
draft-reasoning stage.

\paragraph{Visual understanding.} The engine first generates a detailed,
people-centric description of the image. Short follow-up prompts then ask for
the number of people, the important objects, the actions being performed, the
facial expressions, gestures and body language, the setting, and a description
or role for each person. The answers are parsed into clean concept lists and
stored in a \code{VisualConcepts} object; text utilities normalise the phrases
and extract spatial relationships from the description.

\paragraph{Language and reasoning.} The second sub-stage generates free-form
draft answers for the activity, the relationship and the intention, a visual
rationale that explains them, and an answer to the user's question when one is
provided. Decoding is deterministic: sampling is disabled and a repetition
penalty is applied, so the same input always gives the same output.
Algorithm~\ref{alg:vlm} gives the procedure; \textsc{Generate} denotes greedy
decoding with a repetition penalty.

\begin{algorithm}[htbp]
\caption{Visual understanding and draft reasoning with SmolVLM}
\label{alg:vlm}
\begin{algorithmic}[1]
\Require Preprocessed frame $F$; optional question $Q$
\Ensure Visual concepts $V$; draft reasoning $D$
\State $cap \gets \textsc{Generate}(F, \text{describe-the-people prompt})$
\State $\theta \gets \textsc{ExtractTerms}(cap)$ \Comment{nouns, verbs, adjectives, people}
\ForAll{$k \in \{\text{count, objects, actions, gestures, setting, people}\}$}
  \State $a_k \gets \textsc{Generate}(F, \text{prompt}_k)$
\EndFor
\State $n \gets \textsc{ParseCount}(a_{\text{count}})$; if no number is found, $n \gets \textsc{ParseCount}(cap)$
\State $V.\mathit{objects} \gets \textsc{Split}(a_{\text{objects}}) \cup \theta.\text{nouns}$
\State $V.\mathit{actions} \gets \textsc{Split}(a_{\text{actions}}) \cup \theta.\text{verbs}$
\State $V.\mathit{gestures} \gets \textsc{Split}(a_{\text{gestures}})$;\; $V.\mathit{people} \gets \textsc{Split}(a_{\text{people}}) \cup \theta.\text{people}$
\State $V.\mathit{setting} \gets a_{\text{setting}}$;\; $V.\mathit{spatial} \gets \textsc{SpatialRelations}(cap)$
\If{$n \geq 2$}
  \State $D.\mathit{activity} \gets \textsc{Generate}(F, \text{``what interaction is happening?''})$
  \State $D.\mathit{relationship} \gets \textsc{Generate}(F, \text{``what is their relationship?''})$
\Else
  \State $D.\mathit{activity} \gets \textsc{Generate}(F, \text{``what is the person doing?''})$
  \State $D.\mathit{relationship} \gets \text{``no relationship''}$
\EndIf
\State $D.\mathit{intention} \gets \textsc{Generate}(F, \text{``why are they } D.\mathit{activity}\text{?''})$
\State $D.\mathit{rationale} \gets \textsc{Generate}(F, \text{``what visible evidence shows this?''})$
\If{$Q \neq \textsc{null}$}
  \State $D.\mathit{answer} \gets \textsc{Generate}(F, Q + \text{``answer briefly and give the visual reason''})$
\EndIf
\State \Return $(V, D)$
\end{algorithmic}
\end{algorithm}

\paragraph{Multiple-choice scoring.} For a multiple-choice question the engine
reads the next-token logits for the option letters instead of generating text.
The resulting option probabilities are averaged over several orderings of the
options to reduce letter and position bias (Algorithm~\ref{alg:mcscore}), and
are later used by the reasoning module.

\begin{algorithm}[htbp]
\caption{Debiased multiple-choice scoring}
\label{alg:mcscore}
\begin{algorithmic}[1]
\Require Frame $F$; question $q$; options $o_1,\dots,o_n$; context $c$; number of orderings $k$
\Ensure Probabilities $p_1,\dots,p_n$ with $\sum_i p_i = 1$
\If{$n = 1$}
  \State \Return $(1)$
\EndIf
\State $\Pi \gets \{\text{identity order}\}$; if $k \geq 2$, add the reversed order and $k-2$ cyclic shifts
\State $p \gets (0,\dots,0)$
\ForAll{orderings $\pi \in \Pi$}
  \State $\mathit{prompt} \gets c + q +$ one line ``letter$_i$. $o_{\pi(i)}$'' for each option
  \State append ``Answer with the letter of the best option.'' to $\mathit{prompt}$
  \State $z \gets$ next-token logits of SmolVLM for $(F, \mathit{prompt})$
  \For{$i = 1$ \textbf{to} $n$}
    \State $\ell_i \gets \max_{t \in \text{tokens}(\text{letter}_i)} z_t$
  \EndFor
  \State $s \gets \operatorname{softmax}(\ell)$
  \For{$i = 1$ \textbf{to} $n$}
    \State $p_{\pi(i)} \gets p_{\pi(i)} + s_i$ \Comment{map back to the original option}
  \EndFor
\EndFor
\State \Return $p / |\Pi|$
\end{algorithmic}
\end{algorithm}

\section{Fine-Tuning SmolVLM on VCR}
\label{sec:finetune}

Used zero-shot, SmolVLM sometimes gives reasons that the image does not support.
To adapt it to commonsense reasoning about people, the model is fine-tuned on
the VCR dataset~\cite{zellers2019vcr} with LoRA~\cite{hu2022lora}
(Equation~\ref{eq:lora}). Low-rank adapters of rank 16 are added to the
attention and feed-forward projections of the language model, while the SigLIP
vision encoder and the connector stay frozen, so only a small fraction of the
parameters is trained (approximately \fillin{confirm with
\code{print\_trainable\_parameters()}} trainable parameters).

\paragraph{Data preparation.} The people and objects referred to in each VCR
question are drawn on the image as labelled boxes (\emph{person1},
\emph{person2}, and so on), and the question text uses the same tags. Three
kinds of training example are built from each question: multiple-choice
question answering (\QtoA), multiple-choice selection of the rationale for the
correct answer (\QAtoR), and a free-form answer followed by its visual reason.
The multiple-choice prompts are word for word the prompts used by the option
scoring of Algorithm~\ref{alg:mcscore}, so training directly sharpens the
probabilities that the reasoning module fuses. The order of the options is
shuffled so that no letter is favoured, and the loss is computed only on the
answer tokens.

\paragraph{Training and checkpoint selection.} The model is trained for 20
epochs. After every epoch the adapter is saved and scored on a selection set of
VCR validation questions, and the epoch with the highest \QtoAR{} accuracy is
kept; ties are broken by higher \QtoA{} accuracy and then by the earlier epoch.
The final accuracy is measured on a separate test set of held-out validation
questions, so the choice of epoch does not bias the reported result
(Algorithm~\ref{alg:lora}). In the completed run, epoch 11 was selected. The
selected adapter is merged into the base weights when the model is loaded
through \path{configs/finetuned.yaml}, so inference with the fine-tuned model is
as fast as with the base model.

\paragraph{Implementation.} The procedure is implemented in
\path{finetune/train_vcr_lora.py}. The script streams VCR images from the
dataset's Parquet files, draws the labelled boxes, builds the three kinds of
training example and caches the prepared data. It saves a checkpoint every 50
steps so that an interrupted run can be resumed, and records the training loss,
the per-epoch selection accuracy and the final test accuracy. After the run,
\path{install_adapter.py} copies the selected adapter into the project. The
configuration is listed in Table~\ref{tab:ftconfig}.

\begin{table}[htbp]
\centering
\caption{Fine-tuning configuration}
\label{tab:ftconfig}
\small
\begin{tabularx}{\textwidth}{@{}p{4cm}L@{}}
\toprule
\textbf{Setting} & \textbf{Value}\\
\midrule
Base model & \code{HuggingFaceTB/SmolVLM-Instruct} (2.2B parameters)\\
Adapter & LoRA, rank 16, alpha 32, dropout 0.05\\
Adapted layers & Query, key, value, output, gate, up and down projections of the language model\\
Frozen components & SigLIP vision encoder and connector\\
Training data & Hugging Face VCR dataset (\QtoA, \QAtoR{} and explanation examples)\\
Image resolution & Longest edge 768 pixels\\
Batch size & 2, with gradient accumulation over 8 steps (effective batch of 16)\\
Optimiser & AdamW, learning rate $10^{-4}$, cosine schedule with 5\% warm-up\\
Training length & 20 epochs of 75 steps (1,500 optimiser steps)\\
Checkpoint selection & Best \QtoAR{} on selection questions: epoch 11\\
Test set & 300 held-out VCR validation questions\\
\bottomrule
\end{tabularx}
\end{table}

\begin{algorithm}[htbp]
\caption{LoRA fine-tuning with best-epoch selection}
\label{alg:lora}
\begin{algorithmic}[1]
\Require Base model $M$; training examples $X$; selection set $S_{\text{sel}}$; test set $S_{\text{test}}$; number of epochs $E = 20$
\Ensure Selected adapter $A^*$; test accuracy before and after fine-tuning
\State add LoRA adapters to the language model of $M$; freeze all other weights
\State $r_{\text{base}} \gets \textsc{Evaluate}(M \text{ without adapters}, S_{\text{test}})$
\For{$e = 1$ \textbf{to} $E$}
  \ForAll{mini-batches $B \subset X$}
    \State update the adapters to minimise the cross-entropy of the answer tokens of $B$
  \EndFor
  \State $A_e \gets$ current adapters;\; $r_e \gets \textsc{Evaluate}(M + A_e, S_{\text{sel}})$
\EndFor
\State $e^* \gets \arg\max_e \big(r_e.\text{\QtoAR},\; r_e.\text{\QtoA},\; -e\big)$ \Comment{ties: higher \QtoA, then earlier epoch}
\State $A^* \gets A_{e^*}$;\; $r_{\text{ft}} \gets \textsc{Evaluate}(M + A^*, S_{\text{test}})$
\State \Return $A^*, r_{\text{base}}, r_{\text{ft}}$
\end{algorithmic}
\end{algorithm}

\section{ConceptNet Commonsense Enrichment}
\label{sec:conceptnet}

Visual observations alone may not be enough to distinguish semantically similar
interactions. The enrichment module uses ConceptNet~\cite{speer2017conceptnet}
to add commonsense knowledge to the visual concepts and produces evidence that
is fused into the final prediction. It is divided into a client layer, which
handles backend access, caching, term normalisation and phrase back-off, and an
enrichment layer, which assigns salience, filters evidence and computes label
support.

\paragraph{Key-term extraction.} Candidate terms are collected from the actions,
gestures, roles, objects and setting identified by the VLM, from its description
and draft answers, and from the user's question. Each term receives a salience
weight, highest for interaction-relevant kinds such as actions (1.0) and
gestures (0.9).

\paragraph{Backend.} Three backends are supported: a local SQLite copy of
ConceptNet, the public ConceptNet API with a persistent on-disk cache, and an
offline mode in which taxonomy-based keyword matching remains available. The
local backend is preferred for reproducibility whenever the database is present.
An unknown phrase such as ``shaking hands'' is retried as \code{shake\_hands}
and then as its head word.

\paragraph{Relevance filtering.} For each term the module retrieves its edges
and filters them by relation type and weight. Facts that link the term to the
anchor concepts of the social-interaction labels, or to other concepts observed
in the same image, are ranked highest. Generic hub concepts, which are connected
to almost everything, are suppressed because they carry little discriminative
value.

\paragraph{Label support and ambiguity resolution.} For each activity,
relationship and intention label, the module calculates how strongly the
evidence supports it (Algorithm~\ref{alg:support}). A term contributes support
if it is one of the label's anchor concepts, if a direct edge connects it to an
anchor, if its two-hop neighbourhood overlaps that of the label, or if it
matches one of the label's keywords, weighted by inverse document frequency.
Negative relations such as \emph{Antonym} and \emph{DistinctFrom} reduce the
support. The module therefore produces a numerical support score for every label
together with readable evidence paths, and combines the support with the match
to the VLM draft to produce a preliminary disambiguated label for each task
(Algorithm~\ref{alg:enrich}). The taxonomy, including each label's anchors and
keywords, is kept in \path{configs/taxonomy.yaml} so that it can be changed
without modifying the code.

\begin{algorithm}[htbp]
\caption{ConceptNet commonsense enrichment}
\label{alg:enrich}
\begin{algorithmic}[1]
\Require Visual concepts $V$; draft $D$; optional question $Q$; taxonomy $\mathcal{T}$
\Ensure Facts $\mathcal{F}$; support $S$; draft match $M$; evidence; disambiguated labels $\hat{L}$
\State $E \gets$ terms from $V$, $D$ and $Q$, each with a salience (action 1.0, gesture 0.9, role 0.9, object 0.7, setting 0.6, caption 0.5); keep the top 20
\ForAll{terms $e \in E$}
  \State $(c_e, \mathit{edges}_e) \gets \textsc{Resolve}(e)$ \Comment{phrase $\rightarrow$ lemma $\rightarrow$ head word}
  \If{$c_e$ is shorter than $e$}
    \State $\mathrm{salience}(e) \gets 0.5 \cdot \mathrm{salience}(e)$ \Comment{backed-off sense is less reliable}
  \EndIf
  \State $P_e[v] \gets w_{\text{rel}} \cdot (1 - e^{-\text{weight}})$ for each neighbour $v$ that is not a generic hub
  \State $N_e \gets$ neighbours linked by \emph{Antonym}, \emph{DistinctFrom} or \emph{Not*} relations
\EndFor
\State $\mathcal{F} \gets$ top 6 facts per term, ranked by relevance ($\times 2$ if the neighbour is a label anchor, $\times 1.5$ if it is another observed concept); keep the top 40
\ForAll{tasks $\tau \in \{\text{activity, relationship, intention}\}$}
  \ForAll{labels $L \in \mathcal{T}[\tau]$}
    \State $S_\tau(L) \gets \dfrac{\sum_e w_{e,\tau}\, \textsc{Score}(e, L)}{\sum_e w_{e,\tau}}$ \Comment{Algorithm~\ref{alg:support}}
    \State $M_\tau(L) \gets \max\big(\mathrm{lex}(D_\tau, L),\; 0.8 \max_{e \in \text{draft}} \textsc{Score}(e, L)\big)$
  \EndFor
  \State $S_\tau \gets S_\tau / \max\big(\max_L S_\tau(L),\, 0.3\big)$
  \State $\hat{L}_\tau \gets \arg\max_L \big(S_\tau(L) + M_\tau(L)\big)$
\EndFor
\State \Return $\mathcal{F}, S, M$, evidence paths, $\hat{L}$
\end{algorithmic}
\end{algorithm}

Here $w_{e,\tau}$ is the salience of term $e$ multiplied by the weight of its
kind for task $\tau$; for example, roles count most for the relationship task
and actions most for the activity task.

\begin{algorithm}[htbp]
\caption{Support of one evidence term for one label}
\label{alg:support}
\begin{algorithmic}[1]
\Require Term $e$ with concept $c_e$, neighbour profile $P_e$ and negative neighbours $N_e$; label $L$ with anchors $A_L$, lexicon $X_L$ and neighbour profile $P_L$
\Ensure Support score $s$
\If{$c_e \in A_L$}
  \State \Return $1.0$ \Comment{exact anchor}
\EndIf
\State $s \gets 0$
\ForAll{anchors $a \in A_L$}
  \If{$a \in P_e$}
    \State $s \gets \max(s,\, 0.9\, P_e[a])$ \Comment{direct edge}
  \EndIf
\EndFor
\State $s \gets \max\big(s,\; 0.6 \cdot \min(1,\, 2.5 \cos(P_e, P_L))\big)$ \Comment{two-hop overlap}
\State $H \gets \mathrm{words}(e) \cap X_L$
\If{$H \neq \emptyset$}
  \State $s \gets \max\big(s,\; 0.6 / \min_{w \in H} \mathrm{df}(w)\big)$ \Comment{keyword, IDF-weighted}
\EndIf
\If{$N_e \cap A_L \neq \emptyset$}
  \State $s \gets s - 0.5$ \Comment{contradicting relation}
\EndIf
\State \Return $s$
\end{algorithmic}
\end{algorithm}

\section{Social Interaction Reasoning}
\label{sec:reasoning}

The Social Interaction Reasoning module, \code{SocialInteractionReasoner}, is
the fusion point of Phase I. It combines the visual concepts, the commonsense
evidence, the VLM draft and the optional user question to predict the social
interaction.

\paragraph{Context fusion.} The module builds a structured context from the
scene description, the number of people and their descriptions, the actions,
gestures, objects and setting identified by the VLM, a selection of the most
relevant ConceptNet facts, and the user's question. This context accompanies
the image in every prompt that the module sends to the VLM.

\paragraph{Sequential prediction.} The three tasks are inferred in the order
activity, relationship, intention. The predicted activity is added to the
context for relationship inference, and both the activity and the relationship
are added to the context for intention inference, so that each later prediction
is consistent with the earlier ones.

\paragraph{Candidate shortlisting.} ConceptNet support, VLM draft matching and
lexical evidence are combined into a prior over the taxonomy labels. When the
prior is informative, only the highest-ranked candidates are shortlisted;
otherwise the complete label set is passed to the VLM.

\paragraph{Log-linear fusion.} Let $p_{\text{VLM}}(L)$ be the VLM probability
for label $L$, $p_{\text{CN}}(L)$ the normalised ConceptNet support, and
$m_{\text{draft}}(L)$ the match between the VLM draft and the label. The module
computes
\begin{equation}
  \log P(L) \;\propto\; w_{\text{VLM}} \log p_{\text{VLM}}(L)
  + w_{\text{CN}} \log p_{\text{CN}}(L)
  + w_{\text{draft}}\, m_{\text{draft}}(L)
  \label{eq:fusion}
\end{equation}
with default weights $w_{\text{VLM}} = 1.0$, $w_{\text{CN}} = 0.5$ and
$w_{\text{draft}} = 1.0$. The resulting distribution gives the highest-scoring
label with its confidence, the ranked alternatives with their component
probabilities, and the evidence behind the prediction. A prediction is marked
ambiguous when the probability margin between the top two labels is below 0.15.
Algorithm~\ref{alg:infertask} infers one task in this way.

\begin{algorithm}[htbp]
\caption{Inference of one task with log-linear fusion}
\label{alg:infertask}
\begin{algorithmic}[1]
\Require Task $\tau$; frame $F$; context $c$; visual concepts $V$; support $S_\tau$; draft match $M_\tau$; disambiguated label $\hat{L}_\tau$
\Ensure Prediction for $\tau$: label, confidence, ambiguity flag, alternatives, evidence
\State $\mathcal{L} \gets$ labels of $\tau$ (excluding \code{none\_single} for relationship)
\ForAll{$L \in \mathcal{L}$}
  \State $\mathit{prior}(L) \gets S_\tau(L) + 1.5\, M_\tau(L) + 0.3\, \mathrm{lex}(\text{scene text}, L)$
\EndFor
\If{fewer than 3 labels have $\mathit{prior} > 0.05$}
  \State $C \gets \mathcal{L}$ \Comment{no reliable prior: VLM sees all labels}
\Else
  \State $C \gets$ the 8 labels with the highest prior, always including $\hat{L}_\tau$
\EndIf
\State $p_{\text{VLM}} \gets \textsc{ScoreOptions}(F, \text{question}_\tau, \text{descriptions of } C, c)$ \Comment{Algorithm~\ref{alg:mcscore}}
\If{$\max_{L \in C} S_\tau(L) < 10^{-3}$}
  \State $p_{\text{CN}} \gets$ uniform over $C$
\Else
  \State $p_{\text{CN}} \gets \operatorname{softmax}(S_\tau / T)$, $T = 0.2$
\EndIf
\ForAll{$L \in C$}
  \State $z(L) \gets w_{\text{VLM}} \log p_{\text{VLM}}(L) + w_{\text{CN}} \log p_{\text{CN}}(L) + w_{\text{draft}}\, M_\tau(L)$
\EndFor
\State $P \gets \operatorname{softmax}(z)$; sort $C$ by $P$ in descending order
\State $L^* \gets C_1$;\; $\mathit{ambiguous} \gets \big(P(C_1) - P(C_2) < 0.15\big)$
\State \Return $\big(L^*, P(L^*), \mathit{ambiguous}, C \text{ with } P, p_{\text{VLM}}, p_{\text{CN}}, M_\tau, \text{evidence for } L^*\big)$
\end{algorithmic}
\end{algorithm}

\paragraph{Special cases and explanation.} Scenes with zero or one detected
person are handled explicitly. When only one person is present, the relationship
is set to \code{none\_single}; when no people are detected, the system returns
a no-interaction prediction instead of inferring a social relationship. Once the
three tasks are inferred, an explanation is generated by a prompt that asks the
VLM to relate the prediction to the visible people, gestures, facial
expressions, objects and setting. Algorithm~\ref{alg:reason} applies these steps
to the three tasks in sequence.

\begin{algorithm}[htbp]
\caption{Social interaction reasoning}
\label{alg:reason}
\begin{algorithmic}[1]
\Require Frame $F$; visual concepts $V$; draft $D$; commonsense result $R$; optional question $Q$ and options $O$
\Ensure Social prediction: activity, relationship, intention, explanation, answer
\State $c \gets \textsc{Fuse}(V, R, Q)$ \Comment{caption, concepts, top ConceptNet facts, question}
\If{$V.\mathit{people} = 0$}
  \State \Return no-interaction prediction
\EndIf
\State $A \gets \textsc{InferTask}(\text{activity}, F, c)$ \Comment{Algorithm~\ref{alg:infertask}}
\If{$V.\mathit{people} = 1$}
  \State $\mathit{Rel} \gets$ \code{none\_single} with confidence 1
\Else
  \State $\mathit{Rel} \gets \textsc{InferTask}(\text{relationship}, F, c + \text{``The people are } A\text{.''})$
\EndIf
\State $I \gets \textsc{InferTask}(\text{intention}, F, c + \text{``The people are } A\text{. They are most likely } \mathit{Rel}\text{.''})$
\State $\mathit{pr} \gets$ statement of $(A, \mathit{Rel}, I)$ + top ConceptNet facts + ``explain which visible evidence supports this''
\State $\mathit{Ex} \gets \textsc{Generate}(F, \mathit{pr})$
\If{$Q \neq \textsc{null}$ \textbf{and} $O \neq \textsc{null}$}
  \State $p \gets \textsc{ScoreOptions}(F, Q, O, c)$;\; $\mathit{ans} \gets O_{\arg\max p}$
\ElsIf{$Q \neq \textsc{null}$}
  \State $\mathit{ans} \gets D.\mathit{answer}$
\EndIf
\State \Return $(A, \mathit{Rel}, I, \mathit{Ex}, \mathit{ans})$ with the supporting evidence
\end{algorithmic}
\end{algorithm}

\section{Phase I Integration}

The four modules are integrated through the \code{Phase1Pipeline} class, which
runs them in sequence (Algorithm~\ref{alg:pipeline}). The modules exchange
structured dataclasses rather than free text: \code{PreprocessedFrame} carries
the output of preprocessing, \code{VisualConcepts} and \code{DraftReasoning}
carry the output of the VLM, \code{CommonsenseResult} carries the ConceptNet
evidence, and \code{SocialPrediction} carries the final prediction. Together
these form a \code{Phase1Result}, which can be serialised to JSON with all of
its evidence and module timings.

\begin{algorithm}[htbp]
\caption{Phase I pipeline}
\label{alg:pipeline}
\begin{algorithmic}[1]
\Require Input source $S$, optional question $Q$, optional answer options $O$
\Ensure Phase I result: frame, visual concepts, draft, commonsense evidence, prediction, timings
\State $F \gets \textsc{Preprocess}(S)$ \Comment{Algorithm~\ref{alg:preprocess}}
\State $(V, D) \gets \textsc{SmolVLM}(F, Q)$ \Comment{Algorithm~\ref{alg:vlm}}
\State $R \gets \textsc{EnrichWithConceptNet}(V, D, Q)$ \Comment{Algorithm~\ref{alg:enrich}}
\State $P \gets \textsc{ReasonSocialInteraction}(F, V, D, R, Q, O)$ \Comment{Algorithm~\ref{alg:reason}}
\State \Return $(F, V, D, R, P)$ with the time taken by each stage
\end{algorithmic}
\end{algorithm}

\section{Backend and User Interface}

A FastAPI backend loads the Phase I pipeline once and keeps the model resident
in memory. Each analysis runs as a background job, and the React frontend polls
the job to show which module is currently running. The endpoints are listed in
Appendix~\ref{app:api}.

The interface accepts an uploaded image, a photo captured from the live camera
or an uploaded video, which is analysed at its key moments, together with an
optional question. Results are presented as prediction cards for the activity,
relationship and intention, each with its confidence value and alternatives,
followed by the generated explanation, the supporting ConceptNet evidence, the
visual concepts identified by the VLM and the processing time. The full result
can be downloaded as a JSON report, and a history strip gives access to earlier
analyses from the same session. This interface is an operational demonstration
of Phase I and is distinct from the Phase II explainability dashboard.

\section{Planned Phase II Extension}
\label{sec:phase2plan}

The Phase I data flow already carries what Phase II needs: the normalised tensor
of Equation~\ref{eq:norm} for visual attribution, access to the model through
\code{SmolVLMEngine}, the structured prediction targets, the textual rationale
and the ConceptNet evidence. Phase II will add:

\begin{itemize}[leftmargin=*]
\item \textbf{Visual explanation path:} Grad-CAM~\cite{selvaraju2017gradcam}
  maps for each predicted label, and extraction of the concepts present in the
  most activated regions, giving a visual concept set $\mathcal{V}$.
\item \textbf{Textual explanation path:} extraction of the concepts mentioned in
  the generated explanation, giving a textual concept set $\mathcal{T}$.
\item \textbf{Cross-modal consistency:} the Jaccard overlap of the two sets,
  \begin{equation}
    \mathrm{CMC} = \frac{|\mathcal{T} \cap \mathcal{V}|}{|\mathcal{T} \cup \mathcal{V}|}
    \label{eq:cmc}
  \end{equation}
\item \textbf{Faithfulness evaluation:} the Insertion--Deletion
  metric~\cite{petsiuk2018rise} and the Pointing Game~\cite{zhang2018excitation}.
\item \textbf{Explainability dashboard:} the input frame, the predicted
  interaction, the three labels and the consistency score in a single view.
\end{itemize}

%% =====================================================================
%% CHAPTER 4  EXPERIMENTAL RESULTS AND PERFORMANCE ANALYSIS
%% =====================================================================
\chapter{Experimental Results and Performance Analysis}
\label{ch:results}

This chapter describes the data and the computing environment used, the
experiments conducted in Phase I, and an analysis of their results.

\section{Dataset Description}

\subsection{Visual Commonsense Reasoning (VCR) Dataset}

The VCR dataset~\cite{zellers2019vcr} contains about 290,000 multiple-choice
questions over about 110,000 images taken from movie scenes, drawn from the
LSMDC~\cite{rohrbach2017lsmdc} and MovieClips collections. Each question comes
with four candidate answers and, for the correct answer, four candidate
rationales. Objects and people in each image are detected and referred to in
the text by tags, for example ``Why is [person1] pointing at [person2]?''. The
questions cover explanations, activities, mental states, roles, temporal
inferences and hypothetical situations, which makes them a close match to the
activity, relationship and intention targets of this project. The wrong options
are produced by Adversarial Matching, so they are relevant to the scene and
cannot be rejected on surface cues alone.

This project uses the Hugging Face release of VCR, read directly from its
Parquet files. Since the labels of the official VCR test split are not public,
all evaluation questions are drawn from the official validation split. Three
disjoint sets are used:

\begin{itemize}[leftmargin=*]
\item \textbf{Training set:} questions from the VCR training split, each
  expanded into a \QtoA, a \QAtoR{} and a free-form explanation example. One
  epoch is 75 optimiser steps with an effective batch of 16, that is
  approximately 1,200 training examples per epoch
  (\fillin{number of VCR training questions used}).
\item \textbf{Selection set:} \fillin{N} VCR validation questions, used only to
  choose the epoch.
\item \textbf{Test set:} 300 further VCR validation questions, used neither for
  training nor for choosing the epoch.
\end{itemize}

\subsection{ConceptNet}

ConceptNet 5.5~\cite{speer2017conceptnet} is used as the commonsense knowledge
source. A local SQLite database built from the ConceptNet assertions
(\fillin{English subset; number of edges / size on disk}) allows offline and
reproducible lookups, with the public ConceptNet API and an on-disk cache as a
fallback.

\subsection{Social Interaction Taxonomy}

The label space for the three prediction tasks is defined in
\path{configs/taxonomy.yaml}. Each label has a description, a set of ConceptNet
anchor concepts and a keyword lexicon. Examples of labels are
\emph{Performing}, \emph{Watching together} and \emph{Photographing} for the
activity; \emph{Performer or speaker and audience}, \emph{Teammates or members
of the same group}, \emph{Doctor (or nurse) and patient} and
\code{none\_single} for the relationship; and \emph{To entertain or perform for
others}, \emph{To socialise, relax and spend time together} and \emph{To
comfort, console or emotionally support someone} for the intention. The
taxonomy contains \fillin{a} activity, \fillin{b} relationship and \fillin{c}
intention labels.

\subsection{Test Fixtures and Demonstration Data}

The automated tests use synthetic images, a deterministic stand-in for the VLM
and a small ConceptNet graph, so they run without downloading the model. The
qualitative demonstration uses real photographs, including the concert
photograph of Section~\ref{sec:demo}.

\section{Ecosystem: Hardware and Software Specifications}

\subsection{Hardware}

Fine-tuning was carried out in a Kaggle notebook; the demonstration and the
interactive experiments were run locally. Table~\ref{tab:hardware} lists the
hardware used.

\begin{table}[htbp]
\centering
\caption{Hardware specifications}
\label{tab:hardware}
\small
\begin{tabularx}{\textwidth}{@{}p{3.2cm}LL@{}}
\toprule
\textbf{Component} & \textbf{Fine-tuning (Kaggle)} & \textbf{Inference and demonstration}\\
\midrule
Processor & \fillin{CPU} & \fillin{CPU / SoC}\\
GPU / accelerator & \fillin{e.g. NVIDIA GPU model, memory} & \fillin{CUDA / MPS / CPU}\\
Memory (RAM) & \fillin{GB} & \fillin{GB}\\
Storage & \fillin{GB} & \fillin{GB}\\
Operating system & Linux (Kaggle runtime) & \fillin{OS and version}\\
Training time & \fillin{hours for 20 epochs} & --\\
\bottomrule
\end{tabularx}
\end{table}

\subsection{Software}

The system is organised as a Python package with separate modules for
preprocessing, VLM inference, commonsense enrichment, reasoning, configuration,
data structures and visualisation. Table~\ref{tab:software} lists the main
software components.

\begin{table}[htbp]
\centering
\caption{Software specifications}
\label{tab:software}
\small
\begin{tabularx}{\textwidth}{@{}p{4.2cm}L@{}}
\toprule
\textbf{Component} & \textbf{Technology}\\
\midrule
Programming language & Python \fillin{version}\\
Image/video processing & OpenCV~\cite{bradski2000opencv}, NumPy, PIL\\
Vision-language model & SmolVLM with Hugging Face Transformers~\cite{wolf2020transformers}\\
Deep learning framework & PyTorch~\cite{paszke2019pytorch}\\
Fine-tuning & PEFT (LoRA), Kaggle notebook\\
Natural-language processing & spaCy and project text utilities\\
Commonsense knowledge & ConceptNet, SQLite database and cache\\
Reasoning & Python, NumPy, configurable YAML taxonomy\\
Backend & FastAPI, Uvicorn\\
Frontend & React, Vite, JavaScript\\
Testing & pytest\\
Configuration & YAML\\
\bottomrule
\end{tabularx}
\end{table}

\section{Experiments Conducted}

Five experiments were conducted in Phase I.

\begin{enumerate}[label=\textbf{E\arabic*.},leftmargin=*]
\item \textbf{Functional validation.} The automated test suite was run with
  pytest to check each module and the complete pipeline.
\item \textbf{LoRA fine-tuning with epoch selection.} SmolVLM was fine-tuned
  for 20 epochs with the configuration of Table~\ref{tab:ftconfig}, and the
  adapter was evaluated on the selection set after every epoch.
\item \textbf{Quantitative test evaluation.} The base model and the selected
  fine-tuned model were evaluated on the same 300 held-out test questions on
  \QtoA, \QAtoR{} and \QtoAR.
\item \textbf{Qualitative comparison of explanations.} The free-form answers
  and reasons of the two models were compared for test questions.
\item \textbf{End-to-end demonstration.} The integrated pipeline was run through
  the web interface on real photographs to inspect the predictions, the evidence
  and the processing time.
\end{enumerate}

\section{Performance Analysis}

\subsection{Functional Validation (E1)}

All 23 automated tests passed (Table~\ref{tab:tests}). The preprocessing tests
check image loading, resizing, normalisation, low-light enhancement, blur
detection and frame sampling. The ConceptNet tests check term resolution,
enrichment, label support and access to the local SQLite database, and further
tests cover text normalisation and taxonomy loading. The end-to-end tests run
the complete pipeline and check that commonsense evidence resolves the
prediction when the VLM is uncertain, that scenes with one person or with no
people are handled correctly, and that multiple-choice questions are answered.
The video tests check key-moment selection and summarisation.

\begin{table}[htbp]
\centering
\caption{Automated validation results}
\label{tab:tests}
\small
\begin{tabular}{@{}lcc@{}}
\toprule
\textbf{Test category} & \textbf{Tests} & \textbf{Result}\\
\midrule
Preprocessing & 7 & Passed\\
ConceptNet & 3 & Passed\\
Text and taxonomy & 5 & Passed\\
End-to-end Phase I pipeline & 5 & Passed\\
Video key-moment selection & 3 & Passed\\
\midrule
\textbf{Total} & \textbf{23} & \textbf{23 passed}\\
\bottomrule
\end{tabular}
\end{table}

These tests confirm that the module interfaces work together under controlled
conditions: an input flows through preprocessing, VLM stand-in inference,
ConceptNet enrichment and social reasoning to a structured prediction. Because
the VLM is replaced by a deterministic stand-in, they validate the software, not
the accuracy of the model; accuracy is measured in E2 and E3.

\subsection{Training Convergence (E2)}

The training loss is shown per optimiser step in Figure~\ref{fig:lossstep} and
per epoch on a logarithmic scale in Figure~\ref{fig:lossepoch}. Averaged per
epoch, it falls from 2.03 in the first epoch to 0.08 in epoch 7 and 0.013 in
epoch 11, and it is 0.001 or less from epoch 15 onwards. The loss therefore
converges well before the end of the run.

\begin{figure}[htbp]
\centering
\figimg{0.85\linewidth}{figures/loss_curve.png}
\caption{Training loss of LoRA fine-tuning on VCR over 1,500 optimiser steps (20 epochs)}
\label{fig:lossstep}
\end{figure}

\begin{figure}[htbp]
\centering
\begin{tikzpicture}
\begin{semilogyaxis}[
  width=0.8\linewidth, height=5.6cm,
  xlabel={Epoch}, ylabel={Mean training loss},
  xmin=0.5, xmax=11.5, xtick={1,...,11},
  ymin=0.008, ymax=6,
  ymajorgrids, grid style={gray!20},
  axis line style={gray!60}, tick style={gray!60},
  every axis label/.style={font=\small}, ticklabel style={font=\footnotesize},
]
\addplot[seriesA, line width=1.2pt, mark=*, mark size=1.8pt] coordinates {
  (1,2.029) (2,1.510) (3,1.124) (4,0.684) (5,0.323) (6,0.146)
  (7,0.079) (8,0.050) (9,0.029) (10,0.018) (11,0.013)};
\node[font=\footnotesize, anchor=west] at (axis cs:1.2,2.029) {2.03};
\node[font=\footnotesize, anchor=east] at (axis cs:10.8,0.013) {0.013};
\end{semilogyaxis}
\end{tikzpicture}
\caption[Mean training loss per epoch, epochs 1--11]{Mean training loss per
epoch for epochs 1--11 (logarithmic scale). Epochs 12--20 range from 0.006 to
0.001.}
\label{fig:lossepoch}
\end{figure}

\subsection{Epoch Selection (E2)}

The selection-set accuracy after every epoch is shown in
Figure~\ref{fig:epochacc} and Table~\ref{tab:epochs}. Accuracy rises quickly
over the first three epochs, with \QtoAR{} increasing from 33\% to 46\%, and
then fluctuates between 43\% and 48\% while the training loss keeps falling.
From epoch 12 onwards the accuracy is identical in every epoch, because the
training loss is close to zero and the adapters no longer change.

Epochs 9 and 11 share the highest \QtoAR{} accuracy of 48\%. Epoch 11 is
selected because its \QtoA{} accuracy is higher (63\% against 62\%), and its
\QAtoR{} accuracy of 65\% is the second highest of the run. The final model
therefore uses the epoch-11 adapter rather than the adapter from the last epoch.

\begin{figure}[htbp]
\centering
\begin{tikzpicture}
\begin{axis}[
  width=\linewidth, height=7cm,
  xlabel={Epoch}, ylabel={Selection-set accuracy (\%)},
  xmin=0.5, xmax=20.5, xtick={1,...,20},
  ymin=30, ymax=70,
  ymajorgrids, grid style={gray!20},
  axis line style={gray!60}, tick style={gray!60},
  every axis label/.style={font=\small}, ticklabel style={font=\footnotesize},
  legend style={at={(0.5,1.03)}, anchor=south, legend columns=3,
                draw=none, font=\footnotesize, /tikz/every even column/.append style={column sep=10pt}},
]
\addplot[seriesA, line width=1.2pt, mark=*, mark size=1.8pt] coordinates {
  (1,52) (2,61) (3,63) (4,60) (5,58) (6,60) (7,63) (8,57) (9,62) (10,59) (11,63)
  (12,62) (13,62) (14,62) (15,62) (16,62) (17,62) (18,62) (19,62) (20,62)};
\addlegendentry{\QtoA}
\addplot[seriesB, line width=1.2pt, mark=square*, mark size=1.7pt] coordinates {
  (1,57) (2,59) (3,62) (4,63) (5,63) (6,62) (7,67) (8,62) (9,63) (10,62) (11,65)
  (12,64) (13,64) (14,64) (15,64) (16,64) (17,64) (18,64) (19,64) (20,64)};
\addlegendentry{\QAtoR}
\addplot[seriesC, line width=1.2pt, mark=triangle*, mark size=2.2pt] coordinates {
  (1,33) (2,41) (3,46) (4,43) (5,43) (6,44) (7,47) (8,43) (9,48) (10,44) (11,48)
  (12,47) (13,47) (14,47) (15,47) (16,47) (17,47) (18,47) (19,47) (20,47)};
\addlegendentry{\QtoAR}
\draw[gray, dashed, thick] (axis cs:11,30) -- (axis cs:11,70);
\node[font=\footnotesize, anchor=north west, fill=white, inner sep=1pt] at (axis cs:11.15,39) {selected epoch (11)};
\end{axis}
\end{tikzpicture}
\caption{Selection-set accuracy after each epoch; the dashed line marks the selected epoch (11)}
\label{fig:epochacc}
\end{figure}

\begin{table}[htbp]
\centering
\caption{Per-epoch accuracy on the selection set}
\label{tab:epochs}
\small
\begin{tabular}{@{}ccccc@{}}
\toprule
\textbf{Epoch} & \textbf{Train loss} & \textbf{\QtoA{} (\%)} & \textbf{\QAtoR{} (\%)} & \textbf{\QtoAR{} (\%)}\\
\midrule
1 & 2.029 & 52 & 57 & 33\\
2 & 1.510 & 61 & 59 & 41\\
3 & 1.124 & 63 & 62 & 46\\
4 & 0.684 & 60 & 63 & 43\\
5 & 0.323 & 58 & 63 & 43\\
6 & 0.146 & 60 & 62 & 44\\
7 & 0.079 & 63 & 67 & 47\\
8 & 0.050 & 57 & 62 & 43\\
9 & 0.029 & 62 & 63 & 48\\
10 & 0.018 & 59 & 62 & 44\\
\rowcolor{phaseone}\textbf{11} & \textbf{0.013} & \textbf{63} & \textbf{65} & \textbf{48}\\
12--20 & 0.006 -- 0.001 & 62 & 64 & 47\\
\bottomrule
\end{tabular}
\end{table}

The separation between the two curves -- training loss approaching zero while
the selection accuracy stops improving after epoch 3 -- shows that the later
epochs mostly fit the training examples more closely without generalising
better. Selecting the epoch on held-out data, rather than using the last
checkpoint, guards against this.

\subsection{Test Accuracy (E3)}

The base model and the fine-tuned model (epoch 11) were evaluated on the same
300 held-out validation questions. Table~\ref{tab:testacc} and
Figure~\ref{fig:testbar} give the results. Fine-tuning improves all three tasks
by 14 to 17 percentage points, and the \QtoAR{} accuracy more than doubles, from
15.3\% (46 of 300 questions) to 32.7\% (98 of 300).

\begin{table}[htbp]
\centering
\caption{Accuracy on 300 held-out VCR validation questions}
\label{tab:testacc}
\small
\begin{tabular}{@{}lccc@{}}
\toprule
\textbf{Model} & \textbf{\QtoA{} (\%)} & \textbf{\QAtoR{} (\%)} & \textbf{\QtoAR{} (\%)}\\
\midrule
Chance level & 25.0 & 25.0 & 6.3\\
SmolVLM-Instruct (base) & 44.7 & 41.3 & 15.3\\
SmolVLM + LoRA (epoch 11) & \textbf{58.7} & \textbf{57.7} & \textbf{32.7}\\
\midrule
Improvement (percentage points) & +14.0 & +16.3 & +17.3\\
Correct answers, base $\rightarrow$ fine-tuned & 134 $\rightarrow$ 176 & 124 $\rightarrow$ 173 & 46 $\rightarrow$ 98\\
95\% confidence interval, fine-tuned & $\pm 5.6$ & $\pm 5.6$ & $\pm 5.3$\\
\bottomrule
\end{tabular}
\end{table}

\begin{figure}[htbp]
\centering
\begin{tikzpicture}
\begin{axis}[
  ybar, bar width=16pt, area legend,
  width=0.8\linewidth, height=6.2cm,
  symbolic x coords={qa,qar,qtoar},
  xtick=data, xticklabels={\QtoA, \QAtoR, \QtoAR},
  enlarge x limits=0.25,
  ymin=0, ymax=70, ylabel={Test accuracy (\%)},
  ymajorgrids, grid style={gray!20},
  axis line style={gray!60}, tick style={gray!60},
  every axis label/.style={font=\small}, ticklabel style={font=\footnotesize},
  nodes near coords, nodes near coords style={font=\footnotesize, color=black},
  every node near coord/.append style={/pgf/number format/fixed, /pgf/number format/precision=1},
  legend style={at={(0.5,1.03)}, anchor=south, legend columns=2, draw=none,
                font=\footnotesize, /tikz/every even column/.append style={column sep=10pt}},
]
\addplot[fill=neutralbar, draw=none] coordinates {(qa,44.7) (qar,41.3) (qtoar,15.3)};
\addlegendentry{SmolVLM-Instruct (base)}
\addplot[fill=seriesA, draw=none] coordinates {(qa,58.7) (qar,57.7) (qtoar,32.7)};
\addlegendentry{SmolVLM + LoRA (epoch 11)}
\end{axis}
\end{tikzpicture}
\caption{Test accuracy of the base and fine-tuned models on 300 held-out VCR questions}
\label{fig:testbar}
\end{figure}

\paragraph{Statistical reliability.} With 300 test questions, the 95\% normal
approximation interval for an accuracy $\hat{p}$ is
$\pm 1.96\sqrt{\hat{p}(1-\hat{p})/300}$, which is about $\pm 5.6$ points for
\QtoA{} and \QAtoR{} and $\pm 5.3$ points for \QtoAR{} of the fine-tuned model
(Table~\ref{tab:testacc}). The observed improvements of 14 to 17 points are
well outside these intervals, so the gain from fine-tuning is not an artefact of
the test sample, although the exact accuracy values carry an uncertainty of a
few points.

\paragraph{Selection versus test accuracy.} The test accuracy is lower than the
selection-set accuracy of epoch 11 (63\%, 65\% and 48\%). This is expected: the
selection set is smaller, and the epoch was chosen because it scored well on
that set, so its selection accuracy is optimistically biased. This is why the
result is reported on a separate test set.

\paragraph{Answer and rationale.} If answer and rationale selection were
independent, the expected \QtoAR{} accuracy would be the product of the \QtoA{}
and \QAtoR{} accuracies: $0.587 \times 0.577 \approx 33.9\%$ for the fine-tuned
model, close to the observed 32.7\%, and $0.447 \times 0.413 \approx 18.5\%$ for
the base model, above the observed 15.3\%. Fine-tuning therefore improves both
sub-tasks rather than only one, and the remaining errors are spread across both.

\paragraph{Context from the literature.} For reference, the R2C model reported
65.1\%, 67.3\% and 44.0\% (\QtoA, \QAtoR, \QtoAR) on the full VCR test set, and
human performance was 91.0\%, 93.0\% and 85.0\%~\cite{zellers2019vcr}. These
figures are not directly comparable with ours, which are on a 300-question
sample, and R2C was trained on the full training set with detected object
features. The comparison shows that a 2.2-billion-parameter general-purpose VLM,
adapted with LoRA on a small amount of VCR data, reaches about three quarters of
R2C's joint accuracy while also being able to generate free-form explanations,
and that a large gap to human performance remains.

\subsection{Qualitative Comparison of Explanations (E4)}

Table~\ref{tab:qualitative} compares the free-form explanations of the two
models for four test questions. For the violin and quarrel questions, the
fine-tuned model gives the correct answer together with a visual reason,
whereas the base model only describes what a person is looking at. For the
father--son question, the base model states the relationship correctly while
the fine-tuned model reverses it, although it gives the expected kind of reason
(the age gap). For the dock question, neither model reaches the reference
answer. The fine-tuned answers consistently follow the VCR pattern of an answer
followed by its visual reason.

\begin{table}[htbp]
\centering
\caption{Explanations of the base and fine-tuned models for four test questions}
\label{tab:qualitative}
\footnotesize
\renewcommand{\arraystretch}{1.2}
\begin{tabularx}{\textwidth}{@{}p{1.9cm}L@{}}
\toprule
\multicolumn{2}{@{}l}{\textbf{What is person2 doing?}}\\
Reference & He is playing a violin. He is holding a violin in one hand and applying the bow to it with the other.\\
Base & Looking at the ship.\\
Fine-tuned & He is playing a violin. He is holding a violin and a bow.\\
\midrule
\multicolumn{2}{@{}l}{\textbf{What is happening between person2 and person1?}}\\
Reference & They are quarreling. person2 looks tense and stern, person1 looks surprised.\\
Base & Person2 is looking at person1.\\
Fine-tuned & person1 is telling person2 something she doesn't like. person2 is looking at person1 with a stern expression.\\
\midrule
\multicolumn{2}{@{}l}{\textbf{How do person7 and person1 know each other?}}\\
Reference & person1 is person7's father. They came in together, and person1 is much older than person7.\\
Base & Person7 and person1 know each other. Person7 is the son of person1.\\
Fine-tuned & person7 is person1's father. There is a specific age gap between them, and they are going to embrace.\\
\midrule
\multicolumn{2}{@{}l}{\textbf{Why is person1 and person2 running on the dock?}}\\
Reference & They are trying to board a boat before it leaves. They won tickets on the boat they are trying to reach.\\
Base & Person1 and person2 are running on the dock because they are trying to catch person5.\\
Fine-tuned & person1 and person2 are trying to escape. person1 and person2 are running and person1 is carrying person2 up some stairs.\\
\bottomrule
\end{tabularx}
\end{table}

The father--son example illustrates why Phase II is needed: an explanation can
cite the right kind of evidence and still reach the wrong conclusion. Checking
that the cited evidence corresponds to what the model attends to in the image,
and that the textual and visual explanations agree, is the purpose of the
cross-modal consistency and faithfulness evaluation.

\subsection{End-to-End Demonstration (E5)}
\label{sec:demo}

Figure~\ref{fig:uipred} shows the web interface analysing a photograph of a
singer performing to a crowd at a concert. The system predicts the activity
\emph{Performing} with 93\% confidence, the relationship \emph{Performer or
speaker and audience} with 91\% and the intention \emph{To entertain or perform
for others} with 97\% (Table~\ref{tab:demo}). In each case the next alternative
receives 3\% or less, so none of the predictions is ambiguous. The generated
explanation relates the prediction to visible evidence: the singer's
outstretched arm and microphone, her smile, and the hands of the crowd reaching
towards her.

\begin{figure}[htbp]
\centering
\figimg{\linewidth}{figures/ui_predictions.jpg}
\caption{Web interface: predictions and explanation for a concert photograph}
\label{fig:uipred}
\end{figure}

\begin{table}[htbp]
\centering
\caption{Predictions for the concert photograph}
\label{tab:demo}
\small
\begin{tabularx}{\textwidth}{@{}p{2.2cm}p{4.2cm}cL@{}}
\toprule
\textbf{Task} & \textbf{Prediction} & \textbf{Conf.} & \textbf{Next alternatives}\\
\midrule
Activity & Performing & 93\% & Watching together (3\%), Photographing (1\%)\\
Relationship & Performer or speaker and audience & 91\% & Teammates (3\%), Doctor and patient (2\%)\\
Intention & To entertain or perform for others & 97\% & To socialise (2\%), To comfort someone (1\%)\\
\bottomrule
\end{tabularx}
\end{table}

The supporting evidence and the scene description are shown in
Figure~\ref{fig:uievid}. ConceptNet links the observed concepts to the
predicted labels: ``concert'' and ``sing'' are closely associated with
performing, a performer is typically found on a stage, and a concert is a kind
of show. The VLM identified the setting as a concert, the actions as reaching
out and holding up, the expressions as happy and excited, and the objects as a
water bottle, a microphone and hands. The number of people reached the system's
upper limit of 50, since the crowd is larger than that.

\begin{figure}[htbp]
\centering
\figimg{\linewidth}{figures/ui_evidence.jpg}
\caption{Web interface: supporting ConceptNet knowledge and scene description
for the same photograph}
\label{fig:uievid}
\end{figure}

\paragraph{Processing time.} The complete analysis of the concert photograph
took 33.3 seconds on \fillin{device of Table~\ref{tab:hardware}}. Most of this
time is spent in the VLM: for a scene with two or more people and no question,
Algorithm~\ref{alg:vlm} makes eleven generation calls, the explanation adds one
more, and the option scoring of Algorithm~\ref{alg:mcscore} adds $k$ forward
passes for each of the three tasks. The preprocessing and ConceptNet stages
operate on small arrays and a local database, and contribute comparatively
little. The per-module timings are recorded in every \code{Phase1Result}.

\subsection{Discussion}

Compared with a conventional image-captioning or classification pipeline, the
Phase I system reasons at three levels. At the \emph{visual} level, it
identifies the people, objects, actions, gestures and setting of the scene and
describes them in a caption. At the \emph{commonsense} level, it links these
observations to ConceptNet relations and label-support paths. At the
\emph{decision} level, it fuses the VLM probabilities, the commonsense support
and the draft match into a single prediction. The final prediction is therefore
not a single sentence from the VLM: it retains the intermediate visual concepts,
commonsense facts, component probabilities, alternatives and evidence that the
explainability phase requires.

Table~\ref{tab:status} summarises the status of each Phase I module.

\begin{table}[htbp]
\centering
\caption{Phase I implementation status}
\label{tab:status}
\small
\begin{tabularx}{\textwidth}{@{}p{3.6cm}p{2cm}L@{}}
\toprule
\textbf{Module} & \textbf{Status} & \textbf{Implemented result}\\
\midrule
Input and Preprocessing & Completed & Image/video/live input, quality checks, enhancement, resizing, normalisation and frame sampling\\
Small Vision-Language Model & Completed & SmolVLM integration, visual concept extraction, draft reasoning, deterministic generation and option scoring\\
Fine-tuning (offline) & Completed & LoRA fine-tuning on VCR for 20 epochs; epoch 11 selected; \QtoAR{} on 300 test questions raised from 15.3\% to 32.7\%\\
ConceptNet Commonsense Enrichment & Completed & Concept extraction, ConceptNet lookup, relevance filtering, label support, evidence and ambiguity resolution\\
Social Interaction Reasoning & Completed & Context fusion, activity/relationship/intention inference, log-linear fusion, confidence, alternatives and explanation generation\\
Explanation verification and dashboard & Phase II & Planned for the next phase\\
\bottomrule
\end{tabularx}
\end{table}

\subsection{Limitations}

\begin{itemize}[leftmargin=*]
\item The quantitative evaluation measures the fine-tuned VLM on VCR questions.
  There is not yet a labelled benchmark for the activity, relationship and
  intention predictions of the complete pipeline, so the end-to-end accuracy of
  the fused prediction has been assessed only qualitatively.
\item The fusion weights ($w_{\text{VLM}}$, $w_{\text{CN}}$, $w_{\text{draft}}$)
  and thresholds are set by hand; their effect has not yet been measured by an
  ablation study, and the calibration of the reported confidence values has not
  been measured.
\item The test set has 300 questions, so the reported accuracies carry an
  uncertainty of about $\pm 5$ points.
\item The analysis of one image takes about half a minute, which is acceptable
  for single images but too slow for continuous real-time video.
\item The explanation is generated by the same VLM that makes the prediction,
  and its faithfulness has not yet been verified; this is the subject of
  Phase II.
\end{itemize}

%% =====================================================================
%% CHAPTER 5  SOCIAL IMPACT AND SUSTAINABILITY
%% =====================================================================
\chapter{Social Impact and Sustainability}
\label{ch:impact}

\section{Social Impact}

Systems that interpret social interactions in images are beginning to be used
in everyday tools. Making such systems explainable and lightweight has several
positive applications.

\begin{itemize}[leftmargin=*]
\item \textbf{Assistive technology.} For people who are blind or have low
  vision, a description that goes beyond listing objects -- ``a performer is
  entertaining an audience'', ``two people appear to be arguing'' -- conveys
  the social meaning of a scene. Because the system runs on modest hardware, it
  can be deployed on personal devices without sending images to a remote
  server.
\item \textbf{Support for social-cue learning.} Showing the cues behind an
  interpretation (gestures, expressions, posture, setting) can support people
  who find non-verbal cues hard to read, and can be used in training material
  that teaches how such cues are interpreted.
\item \textbf{Human--robot interaction.} Social and service robots need to
  recognise whether people are busy, in conversation or seeking help before
  approaching them; an inspectable reasoning chain helps developers check that
  such decisions are made for the right reasons.
\item \textbf{Trustworthy AI.} Every prediction comes with its evidence, its
  alternatives and a confidence value, and ambiguous cases are flagged. Users
  can see when the system is uncertain and why it reached a conclusion, which
  supports appropriate rather than blind trust.
\end{itemize}

\section{Ethical Considerations and Risks}

Inferring relationships and intentions from images is sensitive, and the
project recognises the following risks.

\begin{itemize}[leftmargin=*]
\item \textbf{Privacy and identification.} Vision-language models may name
  recognisable people; in the demonstration of Section~\ref{sec:demo}, the
  scene description identifies the performer by name. Identity is not needed
  for the social-interaction task. Prompts and output filtering should prevent
  the system from naming individuals, and face blurring can be applied before
  results are stored.
\item \textbf{Surveillance misuse.} Automatic inference of relationships and
  intentions could be misused to monitor people without consent. The system is
  intended for assistive and research use with consenting users, and processes
  images locally without retaining them by default.
\item \textbf{Bias and stereotyping.} VCR is built from movie scenes, and
  pre-trained VLMs reflect their training data; such sources can encode
  stereotypes about who is a ``doctor'' or a ``patient'', or about gender and
  age roles. Biased performance across demographic groups has been documented
  in commercial vision systems~\cite{buolamwini2018gender}. The visible
  evidence and alternatives produced by this system make such biases easier to
  detect, and a demographic audit of the predictions is planned.
\item \textbf{Over-reliance.} A confident but wrong inference about intention
  can mislead. The ambiguity flag and the alternatives are shown so that users
  treat the output as an interpretation, not a fact.
\end{itemize}

\section{Sustainability}

The design choices of the project reduce the computational and environmental
cost of the system~\cite{strubell2019energy}.

\begin{itemize}[leftmargin=*]
\item \textbf{Small model.} The 2.2-billion-parameter SmolVLM is far smaller
  than the large multimodal models commonly used for visual reasoning, so
  inference runs on a laptop GPU, Apple-silicon MPS or a CPU instead of a
  data-centre accelerator.
\item \textbf{Parameter-efficient training.} LoRA trains only the low-rank
  adapters while the vision encoder, the connector and the language-model
  weights remain frozen. The complete training run was 1,500 optimiser steps on
  a single free cloud notebook (\fillin{GPU-hours}).
\item \textbf{No inference overhead.} The selected adapter is merged into the
  base weights at load time, so the fine-tuned model costs no more to run than
  the base model.
\item \textbf{Reuse instead of retraining.} Commonsense knowledge comes from an
  existing knowledge graph queried from a local database, and the label space is
  a configuration file. New labels can be added without retraining the model.
\item \textbf{Offline operation.} The local ConceptNet database and local model
  make the system usable without a network connection, which reduces both data
  transfer and dependence on paid cloud services.
\item \textbf{Efficient video handling.} Frame sampling skips blurry frames and
  frames that repeat the previous scene, so the model analyses only frames that
  carry new information.
\end{itemize}

\section{Alignment with the Sustainable Development Goals}

The project contributes to the following United Nations Sustainable Development
Goals:

\begin{itemize}[leftmargin=*]
\item \textbf{SDG 3 -- Good Health and Well-being:} assistive scene
  understanding for people with visual impairment and support for social-cue
  learning.
\item \textbf{SDG 9 -- Industry, Innovation and Infrastructure:} affordable,
  explainable AI that runs on commodity hardware.
\item \textbf{SDG 10 -- Reduced Inequalities:} lower hardware requirements widen
  access to multimodal AI, and transparent evidence makes biased predictions
  easier to find and correct.
\end{itemize}

%% =====================================================================
%% CHAPTER 6  CONCLUSIONS AND FUTURE WORK
%% =====================================================================
\chapter{Conclusions and Future Work}
\label{ch:conclusion}

\section{Conclusions}

Phase I of the project has produced a complete, tested prediction pipeline for
explainable visual commonsense reasoning about social interactions. The system
accepts image, video and live-camera input, standardises it with OpenCV,
performs visual understanding and language reasoning with SmolVLM, enriches the
detected concepts with ConceptNet, and fuses the resulting evidence in a
configurable social-interaction reasoning module.

SmolVLM was fine-tuned on the VCR dataset with LoRA for 20 epochs, and the
adapter from epoch 11 was selected on a held-out selection set. On 300 separate
test questions it improves accuracy over the base model from 44.7\% to 58.7\%
for \QtoA, from 41.3\% to 57.7\% for \QAtoR{} and from 15.3\% to 32.7\% for
\QtoAR, gains well beyond the uncertainty of the test sample.

The main output of Phase I is a structured prediction of the activity,
relationship and intention, accompanied by confidence values, alternatives,
evidence and a textual explanation. The modules are connected through a
reusable \code{Phase1Pipeline}, the FastAPI--React interface demonstrates the
complete flow, and all 23 automated tests pass. Phase I objectives 1--10 of
Section~\ref{sec:objectives} have therefore been met, and the implementation
provides a tested foundation for the second phase.

\section{Future Work}

Phase II will extend the system from prediction to verified explanation:

\begin{enumerate}[leftmargin=*]
\item \textbf{Visual explanation path:} Grad-CAM maps over the SmolVLM vision
  encoder for each predicted label, and extraction of the concepts in the most
  activated regions.
\item \textbf{Textual explanation path:} extraction of the concepts cited in the
  generated explanation, independent of the visual path.
\item \textbf{Cross-modal consistency:} Jaccard-based consistency between the
  textual and visual concept sets (Equation~\ref{eq:cmc}).
\item \textbf{Faithfulness evaluation:} the Insertion--Deletion metric and the
  Pointing Game on images with annotated regions.
\item \textbf{Explainability dashboard:} a single view of the input frame, the
  predicted interaction, the three labels, the explanations and the consistency
  score.
\end{enumerate}

In addition, the following improvements are planned:

\begin{itemize}[leftmargin=*]
\item a labelled evaluation set for the activity, relationship and intention
  predictions of the full pipeline, together with an ablation study of the
  fusion weights and of each evidence source;
\item measurement and improvement of confidence calibration;
\item a larger VCR test set and training on more VCR questions;
\item faster inference by batching the VLM prompts, to approach real-time video
  analysis;
\item privacy safeguards that prevent the naming of individuals, and an audit of
  the predictions across demographic groups.
\end{itemize}

%% =====================================================================
%% REFERENCES
%% =====================================================================
\begin{thebibliography}{99}
\setlength{\itemsep}{2pt}

\bibitem{zellers2019vcr}
R.~Zellers, Y.~Bisk, A.~Farhadi, and Y.~Choi, ``From Recognition to Cognition:
Visual Commonsense Reasoning,'' in \emph{Proceedings of the IEEE/CVF Conference
on Computer Vision and Pattern Recognition (CVPR)}, 2019.

\bibitem{lin2019tabvcr}
K.~Lin \emph{et al.}, ``TAB-VCR: Tags and Attributes based VCR Baselines,'' in
\emph{Advances in Neural Information Processing Systems (NeurIPS)}, 2019.

\bibitem{yu2019hgl}
W.~Yu \emph{et al.}, ``Heterogeneous Graph Learning for Visual Commonsense
Reasoning,'' in \emph{Advances in Neural Information Processing Systems
(NeurIPS)}, 2019.

\bibitem{park2020visualcomet}
J.~S. Park, C.~Bhagavatula, R.~Mottaghi, A.~Farhadi, and Y.~Choi,
``VisualCOMET: Reasoning about the Dynamic Context of a Still Image,'' in
\emph{European Conference on Computer Vision (ECCV)}, 2020.

\bibitem{lu2019vilbert}
J.~Lu, D.~Batra, D.~Parikh, and S.~Lee, ``ViLBERT: Pretraining Task-Agnostic
Visiolinguistic Representations for Vision-and-Language Tasks,'' in
\emph{Advances in Neural Information Processing Systems (NeurIPS)}, 2019.

\bibitem{su2020vlbert}
W.~Su \emph{et al.}, ``VL-BERT: Pre-training of Generic Visual-Linguistic
Representations,'' in \emph{International Conference on Learning
Representations (ICLR)}, 2020.

\bibitem{chen2020uniter}
Y.-C. Chen \emph{et al.}, ``UNITER: UNiversal Image-TExt Representation
Learning,'' in \emph{European Conference on Computer Vision (ECCV)}, 2020.

\bibitem{zellers2021merlot}
R.~Zellers \emph{et al.}, ``MERLOT: Multimodal Neural Script Knowledge
Models,'' in \emph{Advances in Neural Information Processing Systems
(NeurIPS)}, 2021.

\bibitem{radford2021clip}
A.~Radford \emph{et al.}, ``Learning Transferable Visual Models From Natural
Language Supervision,'' in \emph{International Conference on Machine Learning
(ICML)}, 2021.

\bibitem{li2023blip2}
J.~Li, D.~Li, S.~Savarese, and S.~Hoi, ``BLIP-2: Bootstrapping Language-Image
Pre-training with Frozen Image Encoders and Large Language Models,'' in
\emph{International Conference on Machine Learning (ICML)}, 2023.

\bibitem{liu2023llava}
H.~Liu, C.~Li, Q.~Wu, and Y.~J. Lee, ``Visual Instruction Tuning,'' in
\emph{Advances in Neural Information Processing Systems (NeurIPS)}, 2023.

\bibitem{marafioti2025smolvlm}
A.~Marafioti \emph{et al.}, ``SmolVLM: Redefining Small and Efficient
Multimodal Models,'' arXiv preprint, 2025.

\bibitem{zhai2023siglip}
X.~Zhai, B.~Mustafa, A.~Kolesnikov, and L.~Beyer, ``Sigmoid Loss for Language
Image Pre-Training,'' in \emph{Proceedings of the IEEE/CVF International
Conference on Computer Vision (ICCV)}, 2023.

\bibitem{speer2017conceptnet}
R.~Speer, J.~Chin, and C.~Havasi, ``ConceptNet 5.5: An Open Multilingual Graph
of General Knowledge,'' in \emph{Proceedings of the AAAI Conference on
Artificial Intelligence}, 2017.

\bibitem{marino2019okvqa}
K.~Marino, M.~Rastegari, A.~Farhadi, and R.~Mottaghi, ``OK-VQA: A Visual
Question Answering Benchmark Requiring External Knowledge,'' in
\emph{Proceedings of the IEEE/CVF Conference on Computer Vision and Pattern
Recognition (CVPR)}, 2019.

\bibitem{marino2021krisp}
K.~Marino, X.~Chen, D.~Parikh, A.~Gupta, and M.~Rohrbach, ``KRISP: Integrating
Implicit and Symbolic Knowledge for Open-Domain Knowledge-Based VQA,'' in
\emph{Proceedings of the IEEE/CVF Conference on Computer Vision and Pattern
Recognition (CVPR)}, 2021.

\bibitem{sun2017domain}
Q.~Sun, B.~Schiele, and M.~Fritz, ``A Domain Based Approach to Social Relation
Recognition,'' in \emph{Proceedings of the IEEE Conference on Computer Vision
and Pattern Recognition (CVPR)}, 2017.

\bibitem{li2017dualglance}
J.~Li, Y.~Wong, Q.~Zhao, and M.~S. Kankanhalli, ``Dual-Glance Model for
Deciphering Social Relationships,'' in \emph{Proceedings of the IEEE
International Conference on Computer Vision (ICCV)}, 2017.

\bibitem{hu2022lora}
E.~J. Hu \emph{et al.}, ``LoRA: Low-Rank Adaptation of Large Language Models,''
in \emph{International Conference on Learning Representations (ICLR)}, 2022.

\bibitem{dettmers2023qlora}
T.~Dettmers, A.~Pagnoni, A.~Holtzman, and L.~Zettlemoyer, ``QLoRA: Efficient
Finetuning of Quantized LLMs,'' in \emph{Advances in Neural Information
Processing Systems (NeurIPS)}, 2023.

\bibitem{selvaraju2017gradcam}
R.~R. Selvaraju \emph{et al.}, ``Grad-CAM: Visual Explanations from Deep
Networks via Gradient-Based Localization,'' in \emph{Proceedings of the IEEE
International Conference on Computer Vision (ICCV)}, 2017.

\bibitem{petsiuk2018rise}
V.~Petsiuk, A.~Das, and K.~Saenko, ``RISE: Randomized Input Sampling for
Explanation of Black-box Models,'' in \emph{British Machine Vision Conference
(BMVC)}, 2018.

\bibitem{zhang2018excitation}
J.~Zhang, S.~A. Bargal, Z.~Lin, J.~Brandt, X.~Shen, and S.~Sclaroff,
``Top-Down Neural Attention by Excitation Backprop,'' \emph{International
Journal of Computer Vision}, vol.~126, pp.~1084--1102, 2018.

\bibitem{kayser2021evil}
M.~Kayser \emph{et al.}, ``e-ViL: A Dataset and Benchmark for Natural Language
Explanations in Vision-Language Tasks,'' in \emph{Proceedings of the IEEE/CVF
International Conference on Computer Vision (ICCV)}, 2021.

\bibitem{zhang2024crossmodal}
Zhang \emph{et al.}, ``Cross-Modal Consistency in Multimodal Large Language
Models,'' 2024.

\bibitem{parcalabescu2025faithfulness}
L.~Parcalabescu and A.~Frank, ``Faithfulness of Vision-Language
Explanations,'' 2024--2025.

\bibitem{bradski2000opencv}
G.~Bradski, ``The OpenCV Library,'' \emph{Dr.\ Dobb's Journal of Software
Tools}, 2000.

\bibitem{zuiderveld1994clahe}
K.~Zuiderveld, ``Contrast Limited Adaptive Histogram Equalization,'' in
\emph{Graphics Gems IV}, P.~S. Heckbert, Ed. Academic Press, 1994,
pp.~474--485.

\bibitem{wolf2020transformers}
T.~Wolf \emph{et al.}, ``Transformers: State-of-the-Art Natural Language
Processing,'' in \emph{Proceedings of the 2020 Conference on Empirical Methods
in Natural Language Processing: System Demonstrations}, 2020.

\bibitem{paszke2019pytorch}
A.~Paszke \emph{et al.}, ``PyTorch: An Imperative Style, High-Performance Deep
Learning Library,'' in \emph{Advances in Neural Information Processing Systems
(NeurIPS)}, 2019.

\bibitem{rohrbach2017lsmdc}
A.~Rohrbach \emph{et al.}, ``Movie Description,'' \emph{International Journal
of Computer Vision}, vol.~123, pp.~94--120, 2017.

\bibitem{buolamwini2018gender}
J.~Buolamwini and T.~Gebru, ``Gender Shades: Intersectional Accuracy
Disparities in Commercial Gender Classification,'' in \emph{Proceedings of the
Conference on Fairness, Accountability and Transparency}, PMLR~81, 2018.

\bibitem{strubell2019energy}
E.~Strubell, A.~Ganesh, and A.~McCallum, ``Energy and Policy Considerations for
Deep Learning in NLP,'' in \emph{Proceedings of the 57th Annual Meeting of the
Association for Computational Linguistics (ACL)}, 2019.

\bibitem{review1}
V.~S. G. \emph{et al.}, ``Explainable Visual Commonsense Reasoning for Social
Interactions -- Review Report I,'' Department of Computer Science and
Engineering, Sri Sivasubramaniya Nadar College of Engineering, 2026.

\end{thebibliography}

%% =====================================================================
%% APPENDICES
%% =====================================================================
\appendix

\chapter{Configuration Parameters}
\label{app:config}

Table~\ref{tab:params} lists the main parameters of the Phase I pipeline, set
in \path{configs/default.yaml} unless stated otherwise.

\begin{table}[htbp]
\centering
\caption{Main configuration parameters of the Phase I pipeline}
\label{tab:params}
\small
\begin{tabularx}{\textwidth}{@{}p{2.6cm}p{5.0cm}L@{}}
\toprule
\textbf{Module} & \textbf{Parameter} & \textbf{Value}\\
\midrule
Preprocessing & Low-light trigger & brightness $< 60$ or contrast $< 25$\\
              & Enhancement & CLAHE on L channel, gamma 0.8\\
              & Blur threshold (Laplacian variance) & $< 60$\\
              & Normalisation $\mu$, $\sigma$ & (0.5, 0.5, 0.5), (0.5, 0.5, 0.5)\\
              & Video key moments (web interface) & up to 7\\
VLM           & Model & \path{HuggingFaceTB/SmolVLM-Instruct}\\
              & Decoding & greedy, repetition penalty\\
              & Fine-tuned model & \path{configs/finetuned.yaml} (epoch-11 adapter, merged)\\
ConceptNet    & Term salience & action 1.0, gesture 0.9, role 0.9, object 0.7, setting 0.6, caption 0.5\\
              & Terms kept & top 20\\
              & Facts kept & top 6 per term, top 40 overall\\
              & Back-off penalty & salience $\times 0.5$\\
              & Contradiction penalty & $-0.5$\\
Reasoning     & Fusion weights & $w_{\text{VLM}} = 1.0$, $w_{\text{CN}} = 0.5$, $w_{\text{draft}} = 1.0$\\
              & Commonsense temperature $T$ & 0.2\\
              & Shortlist size & 8 labels\\
              & Ambiguity margin & 0.15\\
              & Maximum people count & 50\\
\bottomrule
\end{tabularx}
\end{table}

\chapter{API Endpoints}
\label{app:api}

\begin{table}[htbp]
\centering
\caption{FastAPI endpoints implemented for the Phase I demonstration}
\label{tab:api}
\small
\begin{tabularx}{\textwidth}{@{}p{5.2cm}L@{}}
\toprule
\textbf{Endpoint} & \textbf{Purpose}\\
\midrule
\code{GET /api/health} & Reports backend, model and ConceptNet status.\\
\code{POST /api/analyze} & Accepts an image and an optional question and choices, and creates an analysis job.\\
\code{POST /api/analyze-video} & Accepts a video, selects key moments and analyses each of them.\\
\code{GET /api/jobs/\{job\_id\}} & Returns the job status, the active module, the completed stages and the final result.\\
\bottomrule
\end{tabularx}
\end{table}

\end{document}
